%%%%%%%%%%%%%%%%%%%%%%%%%%%%%%%%%%%%%%%%%%%%%%%%%%%%%%%%%%%%%%%%%
\part{Marco teórico}
\label{part:marco-teorico}
%%%%%%%%%%%%%%%%%%%%%%%%%%%%%%%%%%%%%%%%%%%%%%%%%%%%%%%%%%%%%%%%%

Este bloque constituye la base conceptual del trabajo. Su objetivo es presentar, de forma progresiva y autocontenida, todos los fundamentos técnicos que un lector con formación en ingeniería informática necesita para comprender las decisiones de diseño, los experimentos y los resultados que se discuten en las partes posteriores de la memoria. La exposición parte de conceptos elementales (qué es una red neuronal, qué es un token, cómo aprende un modelo) y avanza gradualmente hacia las técnicas específicas sobre las que se articula el TFG, entre ellas la arquitectura transformer, el \textit{fine-tuning} y el PEFT, la destilación de conocimiento, el serving de modelos de lenguaje y la asignación dinámica de peticiones entre modelos de distinto tamaño. Cada concepto se motiva en relación con el problema práctico del proyecto antes de desarrollarse en profundidad, de manera que la teoría no quede aislada del trabajo experimental.

%%%%%%%%%%%%%%%%%%%%%%%%%%%%%%%%%%%%%%%%%%%%%%%%%%%%%%%%%%%%%%%%%
\section{Fundamentos de deep learning}
\label{sec:fundamentos}

Antes de introducir la arquitectura de los modelos de lenguaje conviene fijar una serie de conceptos previos que se utilizarán de forma continua en el resto de la memoria. El objetivo de esta sección no es desarrollar una introducción exhaustiva a todo el campo, sino presentar de forma clara las ideas mínimas necesarias para entender cómo se construyen, entrenan y evalúan los modelos que se estudian en este trabajo. En particular, se introducirán los conceptos de red neuronal, función de pérdida, optimización por gradientes, tokenización y modelo de lenguaje.

\subsection{Redes neuronales artificiales}

Las redes neuronales artificiales constituyen la base de la mayor parte de modelos modernos de \textit{deep learning}. En esencia, una red neuronal es una función parametrizada que transforma una entrada en una salida mediante una secuencia de operaciones matemáticas sencillas. Su interés no reside en la complejidad de cada operación aislada, sino en la capacidad que emerge al combinar muchas de ellas de forma estructurada y al ajustar automáticamente sus parámetros a partir de datos.

La unidad conceptual más simple dentro de una red es la neurona. Una neurona recibe varios valores de entrada, asigna a cada uno de ellos un peso, suma el resultado y añade un sesgo. A continuación, aplica una función de activación que introduce no linealidad en el modelo. De forma simplificada, el cálculo realizado por una neurona puede escribirse como
\[
z = \sum_{i=1}^{n} w_i x_i + b,
\qquad
a = \sigma(z),
\]
donde $x_1, \dots, x_n$ son las entradas, $w_1, \dots, w_n$ los pesos asociados, $b$ es el sesgo, $\sigma$ es la función de activación y $a$ es la salida producida por la neurona. Sin la activación, la red completa se reduciría a una combinación lineal de las entradas y tendría una capacidad expresiva muy limitada.

Cuando varias neuronas se agrupan y procesan conjuntamente la misma entrada, se obtiene una capa. En una capa densa o \textit{fully connected}, cada neurona de una capa recibe información de todas las salidas de la capa anterior. Si una capa recibe un vector de entrada $x \in \mathbb{R}^{n}$ y produce una salida $h \in \mathbb{R}^{m}$, su operación puede expresarse de forma compacta como
\[
h = \sigma(Wx + b),
\]
donde $W \in \mathbb{R}^{m \times n}$ contiene los pesos de la capa y $b \in \mathbb{R}^{m}$ contiene los sesgos. Esta formulación matricial es especialmente importante porque es la forma en la que estas operaciones se implementan de manera eficiente sobre GPU.

\begin{figure}[H]
\centering
\begin{tikzpicture}[
    neuron/.style={circle, draw=black, thick, minimum size=8mm, inner sep=0pt},
    input/.style={neuron, fill=bcdata!30},
    hidden/.style={neuron, fill=bcstudent!25},
    output/.style={neuron, fill=bcteacher!25},
    conn/.style={-Latex, draw=black!60, line width=0.7pt},
    labelbox/.style={rounded corners, draw=black!35, fill=white, align=center, inner sep=5pt}
]

% Entrada
\foreach \i/\y/\txt in {1/1.8/$x_1$,2/0.6/$x_2$,3/-0.6/$x_3$,4/-1.8/$x_n$}{
    \node[input] (I\i) at (0,\y) {};
    \node[left=2mm of I\i] {\small \txt};
}

% Capa oculta 1
\foreach \i/\y in {1/1.2,2/0,3/-1.2}{
    \node[hidden] (H1\i) at (3,\y) {};
}

% Capa oculta 2
\foreach \i/\y in {1/0.8,2/-0.8}{
    \node[hidden] (H2\i) at (6,\y) {};
}

% Salida
\node[output] (O1) at (9,0) {};
\node[right=2mm of O1] {\small $\hat{y}$};

% Conexiones entrada -> oculta 1
\foreach \i in {1,2,3,4}{
    \foreach \j in {1,2,3}{
        \draw[conn] (I\i) -- (H1\j);
    }
}

% Conexiones oculta 1 -> oculta 2
\foreach \i in {1,2,3}{
    \foreach \j in {1,2}{
        \draw[conn] (H1\i) -- (H2\j);
    }
}

% Conexiones oculta 2 -> salida
\foreach \i in {1,2}{
    \draw[conn] (H2\i) -- (O1);
}

% Títulos de capas
\node[above=4mm of I1] {\small \textbf{Capa de entrada}};
\node[above=4mm of H11] {\small \textbf{Capas ocultas}};
\node[above=4mm of O1] {\small \textbf{Salida}};

% Caja explicativa
\node[labelbox] at (4.5,-3.0) {\small En cada neurona\\[2pt] $z=\sum_i w_i x_i + b$\\[2pt] $a=\sigma(z)$};

% Caja lateral
\node[labelbox] at (4.5,3.0) {\small Los pesos $w$ y los sesgos $b$\\son parámetros aprendidos};

\end{tikzpicture}
\caption{Esquema simplificado de una red neuronal densa. Cada neurona combina las señales de la capa anterior mediante pesos, añade un sesgo y aplica una activación no lineal.}
\label{fig:red-neuronal-basica}
\end{figure}

Como se observa en la Figura \ref{fig:red-neuronal-basica}, una red neuronal puede entenderse como una composición de capas. La salida de una capa pasa a ser la entrada de la siguiente, lo que permite construir transformaciones progresivamente más abstractas. Las primeras capas suelen capturar patrones relativamente simples, mientras que las capas posteriores combinan esos patrones para representar relaciones más complejas. Esta idea resulta fundamental para entender por qué las arquitecturas profundas han tenido tanto éxito en tareas de visión, voz y lenguaje.

El término \textit{profunda} hace referencia precisamente a esa acumulación de capas. Una red con pocas capas puede modelar relaciones simples, pero suele necesitar muchas neuronas para aproximar funciones complejas. En cambio, una red profunda distribuye la representación en varios niveles, lo que permite reutilizar patrones intermedios y ganar capacidad expresiva de una forma más eficiente. En la práctica, esta profundidad es uno de los factores que han hecho posible el desarrollo de modelos de gran capacidad como los LLM.

Conviene remarcar que no todas las redes neuronales tienen exactamente la estructura mostrada en la figura. La representación anterior es deliberadamente sencilla y sirve solo como punto de partida conceptual. Los modelos modernos utilizados en procesamiento del lenguaje natural, y en particular los basados en la arquitectura Transformer, no se construyen como una simple pila de capas densas. Sin embargo, siguen respondiendo a la misma idea general. Existen entradas numéricas, parámetros aprendibles, transformaciones lineales, activaciones no lineales y una composición de bloques que produce una salida final.

Finalmente, los pesos y sesgos de todas las capas forman el conjunto de parámetros del modelo. Cuando se afirma que un modelo tiene 7B parámetros, lo que se está diciendo es que contiene aproximadamente siete mil millones de valores numéricos ajustables. Estos parámetros se almacenan en memoria y se modifican durante el entrenamiento para que la red aprenda a resolver la tarea deseada. Este detalle, que puede parecer puramente técnico, es en realidad central en el contexto de este trabajo, ya que el número de parámetros condiciona de forma directa el consumo de memoria, el coste computacional, la latencia de inferencia y, en general, la viabilidad de desplegar el modelo de forma eficiente.
\subsection{Entrenamiento y optimización}

Entrenar una red neuronal consiste en ajustar sus parámetros para que produzca salidas cada vez más cercanas a las deseadas en un conjunto de ejemplos. Dicho de forma simple, el modelo parte de unos parámetros iniciales que no resuelven bien la tarea y, tras ver muchos datos, va corrigiéndolos progresivamente hasta captar patrones útiles. Si denotamos por $\theta$ al conjunto completo de parámetros del modelo, entonces el objetivo del entrenamiento es encontrar unos valores de $\theta$ que minimicen una \textbf{función de pérdida} $\mathcal{L}(\theta)$ sobre un conjunto de datos de entrenamiento $\mathcal{D}$.

La función de pérdida cuantifica el error cometido por el modelo. Cuanto menor es su valor, más próximas están las predicciones del modelo a las respuestas esperadas. Por tanto, entrenar una red neuronal equivale a resolver un problema de optimización en el que se busca el conjunto de parámetros que minimiza dicha pérdida.

En modelos de lenguaje, la función de pérdida más habitual es la \textbf{entropía cruzada} (\textit{cross-entropy}). En este contexto, el modelo recibe un contexto de entrada $x$ y asigna probabilidades a los posibles tokens siguientes. Si el token correcto es $y$, la pérdida asociada a ese ejemplo puede expresarse como
\[
\mathcal{L}(x,y;\theta) = -\log p_\theta(y \mid x).
\]
Esta expresión tiene una interpretación sencilla. Si el modelo asigna una probabilidad alta al token correcto, la pérdida es pequeña. Si le asigna una probabilidad baja, la pérdida aumenta. En consecuencia, minimizar esta función empuja al modelo a concentrar cada vez más probabilidad en las continuaciones correctas.

La minimización de la pérdida se lleva a cabo mediante algoritmos de optimización basados en gradientes. El esquema más básico es el \textbf{descenso por gradiente estocástico} (\textit{stochastic gradient descent}, SGD), aunque en la práctica los LLMs modernos suelen entrenarse con variantes más estables y eficientes, entre las que destaca \textbf{AdamW}. La idea general es siempre la misma. En cada iteración se calcula cómo cambia la pérdida con respecto a cada parámetro del modelo, es decir, su gradiente $\nabla_\theta \mathcal{L}$, y a continuación se actualizan los parámetros en la dirección que reduce dicha pérdida
\[
\theta \leftarrow \theta - \eta \nabla_\theta \mathcal{L},
\]
donde $\eta$ es la \textbf{tasa de aprendizaje} (\textit{learning rate}). Este hiperparámetro controla el tamaño de cada actualización. Si su valor es demasiado alto, el entrenamiento puede volverse inestable. Si es demasiado bajo, el aprendizaje progresa con excesiva lentitud. Por ello, en la práctica se suelen emplear calendarios de aprendizaje en los que la tasa varía a lo largo del entrenamiento, normalmente con una fase inicial de \textit{warmup} seguida de una disminución gradual.

El cálculo eficiente de los gradientes en redes profundas se realiza mediante \textbf{backpropagation}. Este algoritmo aplica de forma sistemática la regla de la cadena a todas las operaciones de la red para obtener la contribución de cada parámetro al error final. Desde un punto de vista práctico, esto implica que cada iteración de entrenamiento se compone de cuatro pasos bien definidos. Primero se ejecuta un \textit{forward pass}, que produce la predicción del modelo. Después se calcula la pérdida. A continuación se realiza el \textit{backward pass}, que obtiene los gradientes. Por último se actualizan los parámetros con el optimizador.

Este ciclo es conceptualmente simple, pero computacionalmente costoso. Durante el entrenamiento no basta con calcular la salida final, ya que también es necesario conservar en memoria muchas activaciones intermedias para poder propagar los gradientes hacia atrás. Por esta razón, entrenar modelos grandes requiere una cantidad elevada de memoria de GPU y un coste computacional muy superior al de la inferencia.

Conviene retener esta idea, ya que será importante en secciones posteriores. Tanto el \textit{finetuning} como la knowledge distillation responden al mismo esquema general de entrenamiento. Lo que cambia entre técnicas no es la lógica básica del proceso, sino la función de pérdida utilizada, los parámetros que se actualizan y la naturaleza de los datos empleados.

\subsection{Tokens y modelos de lenguaje}

Un \textbf{modelo de lenguaje} es un modelo probabilístico que asigna probabilidades a secuencias de texto. Más concretamente, dado un contexto previo, estima qué tokens tienen más probabilidad de aparecer a continuación. Si representamos una secuencia como $(t_1, t_2, \dots, t_n)$, el modelo aprende a estimar la distribución condicional $p(t_{n+1} \mid t_1, \dots, t_n)$ del siguiente token.

Para entender esta formulación, primero hay que aclarar qué es un \textbf{token}. Un token no coincide necesariamente con una palabra completa. En los modelos actuales suele ser una unidad intermedia de texto que puede corresponder a una palabra entera, a una parte de una palabra, a un signo de puntuación o incluso a fragmentos más pequeños. Esta decisión es importante, ya que determina tanto el tamaño del vocabulario como la longitud de las secuencias que procesa el modelo.

Si se trabajara únicamente con palabras completas, el vocabulario sería enorme y aparecería con frecuencia el problema de las palabras fuera de vocabulario. Si se trabajara carácter a carácter, este problema desaparecería, pero las secuencias serían mucho más largas y costosas de procesar. Por ello, los modelos modernos suelen utilizar \textbf{tokenización subword}, que ofrece un equilibrio razonable entre ambas opciones. Entre los métodos más extendidos se encuentra \textit{Byte Pair Encoding} (BPE), que construye un vocabulario de tamaño fijo a partir de fragmentos frecuentes del corpus de entrenamiento. De este modo, las palabras muy comunes pueden representarse con un único token, mientras que las menos frecuentes se descomponen en varias unidades más pequeñas.

Una vez fijada la tokenización, un modelo de lenguaje autoregresivo factoriza la probabilidad conjunta de una secuencia como un producto de probabilidades condicionales
\[
p(t_1, t_2, \dots, t_N) = \prod_{i=1}^{N} p(t_i \mid t_1, \dots, t_{i-1}).
\]
Esta expresión formaliza una idea sencilla. El modelo genera el texto paso a paso, prediciendo en cada instante el siguiente token a partir de todos los anteriores. Por eso se habla de \textbf{generación autoregresiva}. Cada nuevo token pasa a formar parte del contexto y condiciona la predicción del siguiente.

\subsection{Arquitectura transformer}

En las subsecciones anteriores se ha explicado que un modelo de lenguaje recibe una secuencia de tokens y aprende a predecir cuál debería ser el siguiente. Sin embargo, esa formulación solo describe la tarea que resuelve el modelo, no la arquitectura concreta con la que lo hace. Esta distinción es importante, porque un modelo de lenguaje se puede definir por su objetivo de predicción, pero a lo largo del tiempo han existido distintas familias de redes neuronales capaces de abordar ese mismo objetivo.

Antes de la aparición del \textit{transformer}, los modelos de lenguaje neuronales más influyentes se basaban sobre todo en arquitecturas recurrentes, como LSTM y GRU. Estos modelos también procesaban secuencias y también podían utilizarse para predecir el siguiente token, pero lo hacían actualizando un estado oculto paso a paso. Ese enfoque fue útil durante años, aunque imponía limitaciones claras. Por un lado, el entrenamiento era poco paralelizable, ya que cada paso dependía estrictamente del anterior. Por otro, resultaba difícil capturar dependencias largas de forma robusta cuando el contexto crecía.

La aparición del \textit{transformer} cambió este escenario de forma decisiva~\cite{vaswani2017attention}. Su aportación principal no fue cambiar el objetivo del modelo, que seguía siendo modelar secuencias de texto, sino cambiar la manera en que el contexto se representa y se procesa. En lugar de resumir toda la historia previa en un único estado recurrente, el \textit{transformer} permite que cada token consulte directamente otros tokens del contexto mediante mecanismos de atención. Esta idea, unida a una estructura de cómputo mucho más paralelizable, hizo posible entrenar modelos mucho mayores sobre volúmenes de datos mucho más grandes.

Por eso, cuando hoy se habla de LLMs modernos, en la práctica casi siempre se está hablando de modelos basados en la arquitectura \textit{transformer}. En este sentido, el \textit{transformer} no es un detalle interno secundario, sino la pieza que explica tanto la capacidad de estos modelos como buena parte de sus limitaciones operativas. Entenderlo es fundamental no solo para comprender cómo procesa el texto un LLM, sino también para interpretar por qué su inferencia consume tanta memoria, por qué el contexto largo encarece tanto el cálculo y por qué conceptos como \textit{prefill}, \textit{decode}, \textit{KV cache}, TTFT o \textit{throughput} aparecen de forma natural cuando se estudia su despliegue.

A un nivel general, un LLM moderno es una red neuronal basada en la arquitectura \textit{transformer}, entrenada para modelar la distribución condicional del siguiente token dado un contexto previo. Si el contexto viene dado por los tokens $(t_1, t_2, \dots, t_n)$, el modelo estima
\[
p_\theta(t_{n+1} \mid t_1, t_2, \dots, t_n).
\]
Una vez entrenado, generar texto consiste en repetir siempre el mismo procedimiento. El modelo recibe un contexto, calcula una distribución de probabilidad sobre el siguiente token, elige uno según una estrategia de decodificación y lo añade al contexto para continuar. Esta generación paso a paso es lo que se denomina generación autoregresiva.

La cuestión clave es que el \textit{transformer} no resuelve esta tarea propagando la información token a token mediante un estado recurrente, sino construyendo representaciones contextuales en las que cada posición puede incorporar información procedente de otras posiciones relevantes de la secuencia. Esa es la idea central que vertebra toda la arquitectura y la razón por la que el mecanismo de atención ocupa un lugar tan importante en ella.

El punto de partida del modelo sigue siendo el texto tokenizado. Cada token discreto $t_i$ se transforma en un vector denso $e_i \in \mathbb{R}^{d}$ mediante una matriz de embeddings $E \in \mathbb{R}^{|V|\times d}$. Formalmente, el embedding del token puede escribirse como
\[
e_i = E_{t_i,:}.
\]
Estos vectores son aprendidos durante el entrenamiento y constituyen la primera representación continua del texto dentro del modelo. Sin embargo, por sí solos no contienen información sobre el orden en el que aparecen los tokens. Para el modelo, una secuencia de embeddings sin señal posicional adicional sería simplemente un conjunto de vectores, sin una noción clara de qué token va antes o después.

Por esa razón, el \textit{transformer} incorpora información posicional. En el trabajo original de Vaswani et al.~\cite{vaswani2017attention} esta información se añadía mediante codificaciones sinusoidales fijas. En muchos modelos modernos se utilizan \textit{Rotary Position Embeddings} o RoPE~\cite{su2021roformer}. Para esta memoria no es necesario desarrollar toda su formulación matemática. Basta con retener que la representación del token no depende solo de su identidad, sino también de su posición, y que la forma de codificar esa posición influye directamente en la capacidad del modelo para trabajar con contextos largos.

Una vez construida la secuencia de representaciones iniciales, el mecanismo central del \textit{transformer} es la \textbf{self attention}. Su función puede resumirse de forma intuitiva del siguiente modo. Cada token debe decidir qué otros tokens del contexto contienen información útil para representarlo mejor. Para ello, a partir de la entrada de una capa $X \in \mathbb{R}^{n\times d}$, el modelo construye tres proyecciones lineales distintas
\[
Q = XW_Q,\qquad K = XW_K,\qquad V = XW_V.
\]
Aquí $Q$ contiene las \textit{queries}, $K$ las \textit{keys} y $V$ los \textit{values}. La interpretación habitual es la siguiente. Cada token formula una consulta mediante su \textit{query}, cada token ofrece una descripción de lo que puede aportar mediante su \textit{key}, y el contenido que realmente se recuperará si esa relación resulta relevante queda almacenado en el \textit{value}.

La atención se calcula comparando \textit{queries} y \textit{keys} mediante productos escalares. Esas afinidades se normalizan con un \textit{softmax} y se usan después para combinar los \textit{values}
\[
\mathrm{Attn}(Q,K,V)=\mathrm{softmax}\!\left(\frac{QK^\top}{\sqrt{d_k}}\right)V.
\]
Esta fórmula resume una idea muy importante. La salida de cada posición no depende solo de su representación original, sino también de una combinación ponderada de otras posiciones de la secuencia. El factor $\sqrt{d_k}$ evita que los productos escalares crezcan demasiado al aumentar la dimensión interna y saturen el \textit{softmax}.

En un modelo autoregresivo, además, esta atención es causal. Eso significa que la posición $i$ solo puede atender a posiciones anteriores o a sí misma, pero nunca a posiciones futuras. Si no se impusiera esta restricción, el modelo podría ver información que todavía no debería conocer cuando intenta predecir el siguiente token. En la práctica, esta causalidad se implementa mediante una máscara triangular inferior que anula la atención hacia el futuro.

La Figura~\ref{fig:self_attention_transformer} resume esta operación de manera esquemática. Primero, cada representación de entrada se proyecta en \textit{query}, \textit{key} y \textit{value}. Después, las afinidades entre \textit{queries} y \textit{keys} producen una matriz de pesos de atención. Finalmente, esos pesos se utilizan para combinar los \textit{values} y obtener nuevas representaciones contextualizadas.

\begin{figure}[H]
\centering
\begin{tikzpicture}[
font=\footnotesize,
flowarrow/.style={draw,-{Latex[length=2.2mm,width=1.6mm]}},
titlelabel/.style={font=\footnotesize, align=center},
inputbox/.style={draw, rounded corners=2pt, fill=gray!15, minimum width=9mm, minimum height=7mm, align=center},
querybox/.style={draw, rounded corners=2pt, fill=bcstudent!28, minimum width=9mm, minimum height=7mm, align=center},
keybox/.style={draw, rounded corners=2pt, fill=bcteacher!28, minimum width=9mm, minimum height=7mm, align=center},
valuebox/.style={draw, rounded corners=2pt, fill=bcrouter!28, minimum width=9mm, minimum height=7mm, align=center},
outputbox/.style={draw, rounded corners=2pt, fill=bcdata!35, minimum width=9mm, minimum height=7mm, align=center},
matrixcell/.style={draw, minimum width=6mm, minimum height=6mm, inner sep=0pt}
]

\node[titlelabel] at (0.0,3.1) {Entrada};
\node[titlelabel] at (3.6,3.1) {Proyecciones};
\node[titlelabel] at (7.8,3.1) {Pesos de atención};
\node[titlelabel] at (11.5,3.1) {Salida};

\foreach \i/\y in {1/2.2,2/1.45,3/0.7,4/-0.05}{
    \node[inputbox] (x\i) at (0,\y) {$x_{\i}$};
}

\foreach \i/\y in {1/2.2,2/1.45,3/0.7,4/-0.05}{
    \node[querybox] (q\i) at (2.8,\y) {$q_{\i}$};
    \node[keybox]   (k\i) at (3.6,\y) {$k_{\i}$};
    \node[valuebox] (v\i) at (4.4,\y) {$v_{\i}$};
    \draw[flowarrow] (x\i.east) -- (q\i.west);
}

\draw[flowarrow] (4.9,1.0) -- (6.2,1.0);
\node[titlelabel] at (5.55,1.3) {$QK^\top$};

\foreach \x/\y/\c in {
6.9/2.2/bcmeasure!65, 7.55/2.2/gray!12,     8.20/2.2/gray!12,     8.85/2.2/gray!12,
6.9/1.55/bcmeasure!40, 7.55/1.55/bcmeasure!65, 8.20/1.55/gray!12, 8.85/1.55/gray!12,
6.9/0.9/bcmeasure!28,  7.55/0.9/bcmeasure!46,  8.20/0.9/bcmeasure!72, 8.85/0.9/gray!12,
6.9/0.25/bcmeasure!18, 7.55/0.25/bcmeasure!34, 8.20/0.25/bcmeasure!52, 8.85/0.25/bcmeasure!78
}{
    \node[matrixcell, fill=\c] at (\x,\y) {};
}

\node[titlelabel] at (7.9,-0.7) {$\mathrm{softmax}\!\left(QK^\top/\sqrt{d_k}\right)$};

\draw[flowarrow] (9.35,1.0) -- (10.6,1.0);
\node[titlelabel] at (9.95,1.3) {$\cdot V$};

\foreach \i/\y in {1/2.2,2/1.45,3/0.7,4/-0.05}{
    \node[outputbox] (o\i) at (11.5,\y) {$o_{\i}$};
}

\node[titlelabel] at (0.0,-0.7) {$X$};
\node[titlelabel] at (3.6,-0.7) {$Q,\ K,\ V$};
\node[titlelabel] at (11.5,-0.7) {$O$};

\end{tikzpicture}
\caption{Esquema conceptual de la operación de \textit{self attention}. Cada token de entrada se proyecta en \textit{query}, \textit{key} y \textit{value}. A partir de las afinidades entre \textit{queries} y \textit{keys} se obtienen pesos de atención que, en el caso autoregresivo, siguen una estructura causal triangular inferior. Esos pesos se utilizan después para combinar los \textit{values} y producir representaciones contextualizadas. Fuente: elaboración propia.}
\label{fig:self_attention_transformer}
\end{figure}

Hasta aquí la explicación se ha planteado como si existiera una única atención. Sin embargo, los \textit{transformers} reales utilizan \textbf{multi head attention}. La idea es repetir esta operación varias veces en paralelo sobre subespacios distintos de la representación. Si la dimensión oculta es $d$ y se emplean $h$ cabezas, cada una trabaja sobre un subespacio de dimensión $d_k=d/h$. Cada cabeza aprende sus propias matrices de proyección y puede captar patrones diferentes. Después, las salidas de todas las cabezas se concatenan y se proyectan de nuevo al espacio original.

La ventaja de este mecanismo es que el modelo no se ve obligado a resumir todas las relaciones entre tokens en un único patrón de atención. Puede, al menos en principio, especializar distintas cabezas en tipos distintos de dependencia. Algunas pueden centrarse en relaciones locales, otras en patrones de largo alcance y otras en señales más estructurales o semánticas.

Una capa \textit{transformer} completa no contiene solo atención. También incorpora un bloque \textbf{feed forward}, normalmente una MLP de dos transformaciones lineales con una activación no lineal intermedia, como GELU o SwiGLU. Además, tanto el bloque de atención como el bloque \textit{feed forward} están acompañados por conexiones residuales y por una normalización. De forma esquemática, una capa puede escribirse como
\[
X' = X + \mathrm{Attn}(\mathrm{Norm}(X)),\qquad
X'' = X' + \mathrm{FFN}(\mathrm{Norm}(X')).
\]
Esta formulación condensa la lógica general de la capa. La atención mezcla información entre posiciones distintas de la secuencia. El bloque \textit{feed forward} transforma la representación dentro de cada posición de forma independiente. Las conexiones residuales permiten que la información tenga un camino directo a través de la red y facilitan el entrenamiento de modelos profundos. La normalización mantiene controladas las escalas internas y mejora la estabilidad numérica.

La Figura~\ref{fig:transformer_pipeline_clean} resume el flujo funcional del modelo. La secuencia de entrada se transforma primero en embeddings con información posicional. Después, los tokens atraviesan sucesivamente las $L$ capas del modelo, donde se combinan atención y bloque \textit{feed forward}. Finalmente, una cabecera de salida proyecta la representación final al tamaño del vocabulario y produce los logits a partir de los cuales se decide el siguiente token. Durante la inferencia, las claves y valores calculados en atención pueden almacenarse para reutilizarse en pasos posteriores, lo que da lugar al \textit{KV cache}, una pieza fundamental cuando se estudia el coste del \textit{serving}.

\begin{figure}[H]
\centering
\begin{tikzpicture}[
font=\footnotesize,
flowarrow/.style={draw,-{Latex[length=2.2mm,width=1.6mm]}},
blocklabel/.style={font=\footnotesize, align=center},
tokenbox/.style={draw, rounded corners=2pt, minimum width=10mm, minimum height=7mm, align=center},
mainblock/.style={draw, rounded corners=3pt, minimum width=60mm, minimum height=8mm, align=center},
subblock/.style={draw, rounded corners=2pt, minimum width=46mm, minimum height=7mm, align=center},
cachebox/.style={draw, rounded corners=3pt, minimum width=18mm, minimum height=8mm, align=center}
]

\node[tokenbox] (t1) at (0.0,3.4) {$t_1$};
\node[tokenbox] (t2) at (1.4,3.4) {$t_2$};
\node[tokenbox] (t3) at (2.8,3.4) {$t_3$};
\node[tokenbox] (t4) at (4.2,3.4) {$t_4$};
\node[blocklabel] (dots) at (5.5,3.4) {$\dots$};
\node[tokenbox] (tn) at (6.8,3.4) {$t_n$};

\node[mainblock, fill=gray!12] (emb) at (3.4,2.1) {Embeddings + información posicional};

\node[blocklabel] at (3.4,1.35) {$L$ capas transformer};
\node[subblock, fill=bcstudent!22] (attn) at (3.4,0.65) {Atención multi head};
\node[subblock, fill=bcdata!22] (ffn) at (3.4,-0.25) {Bloque feed forward};

\node[mainblock, fill=green!15] (head) at (3.4,-1.55) {LM head\\logits sobre el vocabulario};

\node[tokenbox, fill=yellow!25] (tn1) at (3.4,-2.95) {$t_{n+1}$};

\node[cachebox, fill=orange!18] (cache) at (8.4,0.65) {KV cache};
\node[blocklabel] at (8.4,1.35) {claves y valores};
\node[blocklabel] at (8.0,-2.95) {el token generado\\se añade al contexto};

\foreach \x in {t1,t2,t3,t4,tn}{
    \draw[flowarrow] (\x.south) -- (emb.north);
}
\draw[flowarrow] (emb.south) -- (attn.north);
\draw[flowarrow] (attn.south) -- (ffn.north);
\draw[flowarrow] (ffn.south) -- (head.north);
\draw[flowarrow] (head.south) -- (tn1.north);

\draw[flowarrow, dashed] (attn.east) -- (cache.west);

\end{tikzpicture}
\caption{Esquema funcional simplificado de un \textit{transformer} autoregresivo. La secuencia de entrada se transforma en embeddings con información posicional y atraviesa después un conjunto de capas \textit{transformer}. Cada capa combina atención y bloque \textit{feed forward}. Durante la inferencia, las claves y valores calculados en atención pueden almacenarse en el \textit{KV cache} para reutilizarse en pasos posteriores. La cabecera final produce los logits sobre el vocabulario, a partir de los cuales se decide el siguiente token. Fuente: elaboración propia.}
\label{fig:transformer_pipeline_clean}
\end{figure}

Después de atravesar las $L$ capas del modelo, cada posición queda representada por un vector final $h_i \in \mathbb{R}^{d}$. La cabecera del modelo, o \textit{LM head}, proyecta ese vector al tamaño del vocabulario y produce un conjunto de logits. Al aplicar un \textit{softmax} a esos logits se obtiene una distribución de probabilidad sobre el siguiente token. A partir de esa distribución se decide qué token generar.

La estrategia más simple es la decodificación \textbf{greedy}, que selecciona siempre el token con mayor probabilidad. Otra posibilidad es el \textbf{sampling}, en el que el siguiente token se muestrea de la distribución, quizá modificada con temperatura o truncada mediante esquemas como \textit{top-$k$} o \textit{top-$p$}. En este TFG, la evaluación de calidad se realiza con decodificación greedy, ya que reduce la variabilidad entre ejecuciones y permite comparaciones más limpias entre modelos.

Para cerrar esta subsección conviene retener tres ideas. La primera es que cada token generado tiene que atravesar todas las capas del modelo. La segunda es que la atención y los bloques \textit{feed forward} concentran la mayor parte del coste de inferencia. La tercera es que, cuando el contexto crece o el modelo tiene miles de millones de parámetros, ese coste se vuelve lo bastante alto como para convertir el \textit{serving} en un problema de ingeniería en sí mismo. Esa será precisamente la transición natural hacia las subsecciones siguientes.

\subsection{Fases de la inferencia autoregresiva y KV cache}

La inferencia de un LLM autoregresivo no es una operación uniforme. Aunque desde fuera parezca que el modelo simplemente recibe un prompt y genera una respuesta, internamente el proceso se divide en dos fases con propiedades computacionales muy distintas, \textit{prefill} y \textit{decode}. Distinguir ambas fases es fundamental para entender el comportamiento real de un sistema de \textit{serving}, ya que muchas decisiones de diseño que parecen razonables para una única petición dejan de serlo cuando el servidor debe atender muchas peticiones simultáneas o contextos muy largos.

La fase de \textit{prefill} corresponde al procesamiento inicial del prompt. En esta etapa, todos los tokens de entrada se introducen en el modelo y se procesan conjuntamente. Desde el punto de vista funcional, el objetivo es construir las representaciones internas del contexto completo y calcular, en cada capa, las correspondientes \textit{keys} y \textit{values} asociadas a cada token del prompt. Esta fase aprovecha muy bien la paralelización de la GPU, ya que los tokens del prompt se procesan de forma masiva y los productos matriciales tienen un tamaño suficientemente grande como para mantener una alta ocupación del hardware.

Conviene matizar, sin embargo, que esta paralelización no implica que el coste algorítmico sea lineal con la longitud del prompt. En un transformer con atención causal estándar, la atención dentro del prompt sigue requiriendo comparar posiciones entre sí, de modo que el trabajo asociado a la atención crece de forma cuadrática con la longitud de la secuencia. Aun así, en la práctica el \textit{prefill} suele comportarse como una fase dominada por cómputo, porque la GPU puede explotar de forma muy eficiente ese paralelismo.

Tras completar el \textit{prefill}, el modelo entra en la fase de \textit{decode}. Aquí el patrón de ejecución cambia por completo. En lugar de procesar todos los tokens a la vez, el modelo genera un único token en cada paso, lo añade al contexto y repite el procedimiento hasta alcanzar una condición de parada. Esta naturaleza secuencial introduce una restricción estructural que no puede eliminarse, ya que el token siguiente depende del token recién generado. Por eso, incluso con hardware muy potente, el \textit{decode} mantiene una componente inherentemente serial.

La diferencia esencial entre ambas fases puede resumirse así. En \textit{prefill}, el modelo procesa muchos tokens en paralelo y aprovecha bien la capacidad de cálculo de la GPU. En \textit{decode}, el modelo procesa un solo token nuevo por paso, pero ese token debe atender al contexto acumulado hasta ese momento. Si no existiera ninguna forma de reutilizar resultados previos, en cada paso habría que recalcular las \textit{keys} y \textit{values} de todos los tokens anteriores, lo que volvería prohibitivamente cara la generación.

El mecanismo que evita esa recomputación es el \textbf{KV cache}. La idea es simple. Una vez que un token ya ha sido procesado, sus \textit{keys} y \textit{values} se almacenan en memoria y se reutilizan en los pasos posteriores. De este modo, cuando se genera un nuevo token, solo es necesario calcular sus propias proyecciones y combinar su atención con las \textit{keys} y \textit{values} ya almacenadas del contexto anterior. Esta reutilización no elimina el crecimiento del coste con la longitud del contexto, pero lo reduce de forma drástica y hace viable la inferencia autoregresiva en la práctica.

La consecuencia es que el \textit{decode} deja de estar dominado por la aritmética de la atención y pasa a estar dominado, en gran medida, por el acceso a memoria. En cada paso hay que leer repetidamente pesos del modelo y, además, recorrer el KV cache acumulado hasta ese momento. Por eso suele decirse que el \textit{prefill} es una fase \textit{compute-bound}, mientras que el \textit{decode} es una fase \textit{memory-bound}. Esta asimetría es una de las ideas más importantes de todo el problema de \textit{serving}, ya que explica por qué un sistema puede mostrar una utilización alta de GPU durante el procesamiento del prompt y, sin embargo, bajar mucho su eficiencia durante la generación token a token.

La Tabla~\ref{tab:prefill-decode-kv} resume las diferencias más relevantes entre ambas fases.

\begin{table}[H]
\centering
\small
\begin{tabularx}{0.98\textwidth}{>{\raggedright\arraybackslash}p{3.1cm} >{\raggedright\arraybackslash}X >{\raggedright\arraybackslash}X}
\toprule
\textbf{Aspecto} & \textbf{Prefill} & \textbf{Decode} \\
\midrule
Unidad de trabajo dominante &
Se procesa el prompt completo de una vez &
Se genera un único token nuevo en cada paso \\

Paralelización &
Muy alta, ya que muchos tokens se evalúan simultáneamente &
Baja, porque la generación es secuencial por construcción \\

Cuello de botella principal &
Cómputo, especialmente productos matriciales grandes &
Memoria, sobre todo lectura repetida de pesos y del KV cache \\

Dependencia respecto al contexto &
La atención sobre el prompt crece con la longitud de la secuencia &
Cada paso depende del contexto acumulado y su coste crece con ese contexto \\

Métrica de latencia más afectada &
Suele dominar el \textit{time to first token} cuando el prompt es largo &
Suele dominar el tiempo por token generado y la latencia total de respuestas largas \\

Sensibilidad al batching &
Se beneficia mucho del procesamiento en lote &
Puede degradarse antes por presión sobre memoria y por interferencia entre secuencias \\

\bottomrule
\end{tabularx}
\caption{Comparación funcional entre las fases de \textit{prefill} y \textit{decode} en la inferencia autoregresiva de un LLM. Fuente: elaboración propia.}
\label{tab:prefill-decode-kv}
\end{table}

Desde el punto de vista de las métricas de sistema, esta separación también es muy útil. El \textit{time to first token} depende en gran medida del coste de procesar el prompt y de completar el primer paso de generación. Por eso, cuando el contexto de entrada es largo, el \textit{prefill} suele pesar mucho en esta métrica. En cambio, una vez que la respuesta ya ha comenzado, el ritmo al que se generan los tokens siguientes depende sobre todo del \textit{decode}. Esto afecta directamente al \textit{throughput} efectivo y a la percepción de fluidez del sistema.

Otro punto clave es el tamaño del KV cache. Su crecimiento es especialmente importante porque consume memoria de forma proporcional no solo a la longitud del contexto, sino también al número de secuencias concurrentes. Una forma útil de expresar esta memoria es la siguiente
\[
M_{\mathrm{KV}} \approx B \, T \, L \, H_{\mathrm{KV}} \, d_{\mathrm{head}} \, 2 \, s,
\]
donde $B$ es el número de secuencias activas, $T$ el número de tokens almacenados por secuencia, $L$ el número de capas, $H_{\mathrm{KV}}$ el número de cabezas de KV, $d_{\mathrm{head}}$ la dimensión por cabeza, el factor $2$ refleja que se almacenan tanto \textit{keys} como \textit{values}, y $s$ es el número de bytes por elemento, que en FP16 vale 2.

Si se quiere expresar la memoria por token y por secuencia, basta con fijar $B=1$ y considerar un único token adicional
\[
M_{\mathrm{KV,token}} \approx L \, H_{\mathrm{KV}} \, d_{\mathrm{head}} \, 2 \, s.
\]
Esta expresión es útil porque permite estimar rápidamente el impacto de aumentar el contexto o el número de peticiones concurrentes.

En una configuración genérica sin \textit{grouped-query attention}, por ejemplo un modelo de alrededor de 7B parámetros con 32 capas, 32 cabezas de KV, dimensión de cabeza 128 y precisión FP16, el coste es
\[
M_{\mathrm{KV,token}} \approx 32 \cdot 32 \cdot 128 \cdot 2 \cdot 2
= 524\,288\ \text{bytes},
\]
es decir, aproximadamente 0,5 MiB por token y por secuencia. Con este orden de magnitud, un contexto de 4\,000 tokens puede consumir alrededor de 2 GiB solo en KV cache, mientras que un contexto de 32\,000 tokens puede acercarse a 16 GiB. Estas cifras ayudan a entender por qué el contexto largo se convierte tan rápidamente en un problema de memoria incluso en GPUs con gran capacidad.

Conviene añadir una precisión importante. En muchos modelos actuales no se almacenan tantas cabezas de KV como cabezas de atención total, porque se emplean variantes como \textit{grouped-query attention} o \textit{multi-query attention}. En esos casos, el valor real de $H_{\mathrm{KV}}$ es menor y el consumo de KV cache se reduce de forma apreciable. Aun así, la dependencia lineal con la longitud del contexto y con el número de secuencias concurrentes sigue presente, de modo que el problema no desaparece, solo se atenúa.

En la práctica, motores de \textit{serving} como vLLM tratan de gestionar esta presión de memoria con técnicas específicas de planificación y de organización del cache, entre ellas \textit{paged attention}. La idea general es reducir fragmentación, reutilizar mejor la memoria disponible y permitir que múltiples secuencias convivan de forma más eficiente. Sin embargo, estas técnicas no cambian la estructura básica del problema. El \textit{prefill} sigue estando asociado a una fase más favorable al cómputo intensivo y el \textit{decode} sigue estando limitado en gran parte por el movimiento de datos y por el tamaño del contexto acumulado.

Para este TFG, esta distinción tiene una consecuencia directa. No basta con comparar modelos solo por calidad de respuesta. También es necesario entender cómo cambia su comportamiento durante el \textit{serving}, ya que un modelo pequeño no solo reduce el coste de los pesos, sino que también puede modificar la latencia, el \textit{throughput} y la presión sobre memoria cuando se sirve a múltiples usuarios o con prompts largos. Esa es precisamente una de las razones por las que el análisis conjunto de calidad y eficiencia resulta central en el trabajo.

\subsection{Métricas de servicio e indicadores de memoria}

Para evaluar el coste de inferencia no basta con medir si un modelo responde bien o mal. También es necesario caracterizar cómo se comporta el sistema cuando sirve peticiones reales. En este TFG se emplean cuatro métricas de servicio estándar, TTFT, TPOT, latencia total y throughput. Junto a ellas se reporta además la memoria GPU como indicador complementario, ya que condiciona de forma directa la concurrencia admisible y la longitud máxima de contexto que el sistema puede sostener.

La primera métrica es el \textbf{TTFT} (\textit{Time To First Token}), que mide el tiempo transcurrido entre el inicio de la petición y la recepción del primer token de salida. Esta definición debe interpretarse con cuidado, porque su valor exacto depende del punto donde se instrumenta la medida. Si la medición se hace extremo a extremo desde el cliente, el TTFT incluye la latencia de red, la posible espera en cola dentro del motor de inferencia y el coste asociado al \textit{prefill}. Si, en cambio, la medición se realiza solo en el servidor, las componentes de red quedan fuera. En cualquier caso, es la métrica que más influye en la percepción de responsividad de una aplicación interactiva.

La segunda métrica es el \textbf{TPOT} (\textit{Time Per Output Token}), que cuantifica el tiempo medio necesario para generar cada token de salida una vez emitido el primero. Esta métrica refleja el comportamiento del sistema durante la fase de \textit{decode} y, por tanto, está muy relacionada con el acceso a memoria, el tamaño del batch efectivo y el grado de contención entre peticiones concurrentes. En muchos trabajos también se reporta su inverso, expresado en tokens por segundo y por petición, porque esa forma resulta más intuitiva al comparar configuraciones.

La tercera métrica es la \textbf{latencia total}, que mide la duración completa de una petición desde su inicio hasta la emisión del último token. Si se denota por $n_{\text{out}}$ el número de tokens generados, una aproximación útil es
\[
\ell_{\text{total}} \;\approx\; \text{TTFT} + (n_{\text{out}} - 1)\cdot \text{TPOT}.
\]
Esta expresión no pretende modelar todos los detalles internos del sistema, pero sí resume bien la estructura general del proceso. Primero hay un coste inicial alto, asociado sobre todo al \textit{prefill} y a la obtención del primer token. Después aparece un coste incremental por cada token adicional generado. En la práctica, la latencia total depende tanto de la longitud del prompt como de la longitud de la respuesta.

La cuarta métrica es el \textbf{throughput}, que mide la cantidad agregada de tokens servidos por segundo a nivel de sistema. Esta métrica no debe confundirse con la velocidad de generación de una petición individual. Un sistema puede aumentar su throughput total atendiendo más secuencias en paralelo, aunque eso empeore el TTFT o el TPOT de cada petición concreta. Aparece así el compromiso clásico entre latencia y capacidad agregada del sistema, que es uno de los ejes centrales en cualquier estudio de \textit{serving}.

Además de estas cuatro métricas de servicio, en este TFG se reporta también la \textbf{memoria GPU} ocupada. No se trata de una métrica de latencia en sentido estricto, pero sí de un indicador operativo esencial. La memoria consumida por los pesos del modelo, el KV cache y las activaciones intermedias determina cuántas peticiones pueden coexistir, qué longitud de contexto puede admitirse y si es posible servir varios modelos en paralelo dentro de la misma GPU. Por este motivo, la memoria se considera aquí una variable de soporte indispensable para interpretar correctamente el resto de resultados.

La Tabla~\ref{tab:metricas-servicio} resume estas métricas e indicadores y su papel dentro del TFG.

\begin{table}[H]
\centering
\caption{Métricas de servicio e indicador complementario de memoria utilizados en este TFG. Fuente: elaboración propia a partir de~\cite{vllm2023, vllm-docs}.}
\label{tab:metricas-servicio}
\footnotesize
\begin{tabularx}{\textwidth}{>{\raggedright\arraybackslash}p{2.6cm} >{\raggedright\arraybackslash}X >{\raggedright\arraybackslash}p{4.6cm}}
\toprule
\textbf{Magnitud} & \textbf{Qué mide} & \textbf{Interpretación en el TFG} \\
\midrule
TTFT & Tiempo hasta el primer token & Mide la responsividad inicial del sistema y está muy influido por el \textit{prefill}, la cola interna y, si la medida es extremo a extremo, también por la red. \\
\addlinespace
TPOT & Tiempo medio por token generado tras el primero & Refleja el comportamiento del \textit{decode} en régimen estable y depende en gran medida del acceso a memoria y del batch efectivo. \\
\addlinespace
Latencia total & Duración completa de la petición & Resume el tiempo total de respuesta y combina el coste inicial del primer token con el coste incremental del resto de tokens generados. \\
\addlinespace
Throughput & Tokens servidos por segundo a nivel de sistema & Mide la capacidad agregada del motor de inferencia y permite estudiar el compromiso entre rendimiento global y latencia por petición. \\
\addlinespace
Memoria GPU & VRAM ocupada por pesos, KV cache y activaciones & No es una métrica de servicio en sentido estricto, pero condiciona la concurrencia, la longitud de contexto y la viabilidad del despliegue. \\
\bottomrule
\end{tabularx}
\end{table}

A la hora de reportar resultados, no todas estas magnitudes deben resumirse del mismo modo. En TTFT y latencia total resulta más informativo mostrar percentiles, especialmente p50 y p95, porque la media no describe bien el comportamiento de la cola de la distribución cuando existe variabilidad entre peticiones o episodios de congestión. En cambio, para throughput suele ser más natural reportar la media sobre una ventana de medida estable, ya que en este caso interesa captar la capacidad agregada del sistema más que la variabilidad entre observaciones individuales.

Estas métricas forman la base de la evaluación experimental que se presenta más adelante. Gracias a ellas es posible comparar modelos no solo por calidad, sino también por coste operativo real. Esta doble perspectiva es imprescindible en un trabajo como el presente, donde el objetivo no es únicamente obtener buenas respuestas, sino estudiar el equilibrio entre calidad, latencia, capacidad de servicio y consumo de recursos en un entorno de supercomputación.

\section{Optimización del modelo. Panorama de técnicas, fine tuning y LoRA}
\label{sec:optim-modelo}

\subsection{Panorama de técnicas para reducir coste y adaptar LLMs}
\label{sec:panorama-opt}

Antes de entrar en detalle en \textit{knowledge distillation}, conviene situarla dentro de un mapa más amplio de técnicas que hoy se utilizan para hacer más eficiente el uso de LLMs. Esta contextualización es importante por dos motivos. En primer lugar, permite entender mejor qué problema resuelve exactamente cada familia de métodos. En segundo lugar, deja claro que la destilación no compite necesariamente con el resto de técnicas, sino que suele combinarse con ellas en sistemas reales.

Cuando se habla de hacer un LLM más barato o más fácil de desplegar, en realidad se están mezclando varios objetivos distintos. A veces se quiere reducir memoria. Otras veces se quiere disminuir latencia. En otros casos lo importante es adaptar el modelo a un dominio concreto sin asumir el coste de reentrenar todos sus parámetros. También hay situaciones en las que el modelo no cambia en absoluto y lo que se optimiza es únicamente la forma de ejecutarlo. Por eso resulta útil distinguir entre técnicas que comprimen el modelo, técnicas que lo adaptan y técnicas que optimizan el \textit{serving}.

Dentro de las técnicas que actúan directamente sobre el modelo, la primera familia relevante es la \textbf{cuantización}. Su idea es reducir la precisión numérica de los pesos, y en algunos casos también de las activaciones, pasando de formatos como FP16 a INT8, INT4 o configuraciones mixtas. El efecto inmediato es una reducción del consumo de memoria y, cuando el hardware lo permite, también una mejora del rendimiento de inferencia. La ventaja principal de la cuantización es que puede aplicarse sin reentrenamiento en su variante \textit{post training}. A cambio, introduce una posible pérdida de calidad que depende del modelo, del esquema de cuantización y del nivel de compresión empleado.

Otra familia importante es el \textbf{pruning}. En este caso no se cambia la precisión de los parámetros, sino que se eliminan parámetros o estructuras internas que se consideran poco relevantes. El pruning puede ser no estructurado, si elimina conexiones individuales, o estructurado, si elimina bloques completos, neuronas, cabezas de atención o canales. Desde el punto de vista práctico, el pruning estructurado suele ser más útil para despliegue, porque produce modelos que pueden aprovechar mejor hardware estándar sin depender tanto de kernels especializados para sparsity. En general, después del pruning suele ser necesario un ajuste adicional para recuperar parte de la calidad perdida.

La tercera gran familia es la \textbf{destilación de modelos} o \textit{knowledge distillation}. Esta técnica ocupa el lugar central en el TFG y se desarrolla en detalle en el apartado siguiente, por lo que aquí solo interesa situarla dentro del panorama general. A diferencia de cuantización y pruning, que transforman el mismo modelo original, la destilación entrena un modelo nuevo y más pequeño para aproximar el comportamiento de otro más grande. Esa diferencia es importante. La destilación suele requerir más coste de entrenamiento que una cuantización \textit{post training}, pero ofrece una flexibilidad mayor, porque permite elegir el tamaño del \textit{student}, cambiar de arquitectura e incluso construir varios modelos pequeños a partir de un mismo \textit{teacher}.

Una segunda familia de técnicas no busca comprimir el modelo en inferencia, sino \textbf{adaptarlo} con un coste de entrenamiento mucho menor. Aquí se sitúan los métodos de \textit{parameter efficient fine tuning}, habitualmente agrupados bajo el acrónimo \textbf{PEFT}. Dentro de esta familia se encuentran técnicas como \textit{adapters}, \textit{prompt tuning}, \textit{prefix tuning} y LoRA. Su idea común es que no siempre es necesario actualizar todos los parámetros del modelo para especializarlo. En muchos casos basta con modificar una fracción muy pequeña o con añadir una corrección entrenable de tamaño reducido. Estas técnicas son relevantes para este TFG porque permiten explorar adaptaciones de bajo coste y porque pueden combinarse con destilación en distintas fases del pipeline.

Una tercera línea distinta son las arquitecturas con \textbf{activación dispersa}, entre ellas \textit{mixture of experts}. En estos modelos no todos los parámetros participan en el procesamiento de cada token. Aunque el número total de parámetros puede ser muy alto, el coste efectivo por token es menor porque solo se activa un subconjunto de expertos. No se trata exactamente de una técnica de compresión aplicada a posteriori, sino de una decisión arquitectónica de diseño. Su inclusión en este panorama es útil porque muestra que la eficiencia no depende únicamente del tamaño nominal del modelo, sino también de cuántos parámetros intervienen realmente durante la inferencia.

Por último, existe una capa adicional de optimización que no modifica el modelo, sino el modo en que se ejecuta. Aquí se sitúan técnicas de \textbf{optimización de serving} como \textit{continuous batching}, \textit{paged attention}, \textit{flash attention} o \textit{speculative decoding}. Todas ellas se desarrollan más adelante en la parte dedicada al \textit{serving}. En este punto basta con remarcar que son técnicas ortogonales a la compresión y a la adaptación. Un modelo cuantizado, podado o destilado sigue pudiendo beneficiarse de estas optimizaciones en producción.

La Tabla~\ref{tab:panorama-optimizacion} resume este mapa general y sitúa cada familia dentro del problema que aborda.

\begin{table}[H]
\centering
\caption{Panorama general de técnicas de optimización, adaptación y serving de LLMs. Fuente: elaboración propia.}
\label{tab:panorama-optimizacion}
\footnotesize
\begin{tabularx}{\textwidth}{>{\raggedright\arraybackslash}p{2.8cm} >{\raggedright\arraybackslash}p{2.9cm} >{\raggedright\arraybackslash}X >{\raggedright\arraybackslash}p{3.5cm}}
\toprule
\textbf{Familia} & \textbf{Qué modifica} & \textbf{Objetivo principal} & \textbf{Relación con el TFG} \\
\midrule
Cuantización & Precisión numérica de pesos y, a veces, activaciones & Reducir memoria y acelerar inferencia con coste de aplicación bajo & Compatible con los modelos destilados obtenidos en el proyecto \\
\addlinespace
Pruning & Parámetros o estructuras internas del modelo & Eliminar redundancia y reducir tamaño efectivo del modelo & No se explora directamente, pero es complementario a la destilación \\
\addlinespace
Knowledge distillation & Entrena un modelo nuevo a partir de otro mayor & Obtener un \textit{student} más barato manteniendo parte de la calidad del \textit{teacher} & Técnica central del TFG \\
\addlinespace
PEFT y LoRA & Solo una pequeña parte del espacio de parámetros & Adaptar un modelo con muy poco coste de entrenamiento & Se utiliza como vía eficiente de experimentación y comparación \\
\addlinespace
Arquitecturas dispersas & La arquitectura y el patrón de activación & Reducir el coste efectivo por token sin reducir necesariamente los parámetros totales & Se mencionan como contexto, pero quedan fuera del alcance \\
\addlinespace
Optimización de serving & La ejecución del modelo, no sus pesos & Mejorar latencia, throughput y uso de memoria en producción & Se combina con los modelos del TFG mediante vLLM \\
\bottomrule
\end{tabularx}
\end{table}

Con este panorama en mente, la elección de \textit{knowledge distillation} como pieza central del TFG se entiende mejor. En primer lugar, produce un modelo nuevo y más pequeño, lo que permite servirlo en paralelo al modelo grande y construir cascadas o estrategias de \textit{routing}. En segundo lugar, no cierra la puerta a técnicas adicionales. Un \textit{student} destilado puede cuantizarse después, puede adaptarse con LoRA y puede servirse con motores optimizados. En tercer lugar, permite estudiar de forma especialmente clara el equilibrio entre calidad y coste, que es precisamente el núcleo experimental del trabajo.

\subsection{Fine tuning y adaptación eficiente}

Aunque la técnica principal del TFG es la destilación de conocimiento, en la práctica muchos experimentos de destilación se implementan como una forma particular de \textit{fine tuning} sobre el modelo pequeño. Por este motivo, antes de entrar en \textit{knowledge distillation}, conviene entender con claridad qué significa ajustar un LLM, qué parámetros se modifican, qué coste tiene hacerlo y por qué el \textit{full fine tuning} tiene un papel tan importante en la línea experimental final del trabajo.

Un LLM no llega al \textit{fine tuning} como una red vacía. Llega después de una fase de preentrenamiento masivo, en la que ha aprendido regularidades generales del lenguaje, patrones de razonamiento, estructuras gramaticales y cierta información factual. Ese conocimiento queda codificado en sus parámetros, que se pueden denotar como $\theta_{\mathrm{pre}}$. El objetivo del \textit{fine tuning} no es volver a aprender todo desde cero, sino desplazar esos parámetros hacia una zona más adecuada para una tarea, un dominio o un formato de respuesta concretos.

Desde el punto de vista matemático, el \textit{fine tuning} sigue siendo un problema de optimización basado en gradientes. Dado un conjunto de datos $\mathcal{D} = \{(x_i,y_i)\}$ y una función de pérdida $\mathcal{L}$, el objetivo puede escribirse como
\[
\theta^\star
=
\arg\min_{\theta}
\mathbb{E}_{(x,y)\sim \mathcal{D}}
\big[
\mathcal{L}(f_\theta(x),y)
\big],
\]
pero con una diferencia fundamental respecto al entrenamiento desde cero. La optimización parte de $\theta_{\mathrm{pre}}$, no de una inicialización aleatoria. Esto hace que el modelo ya disponga de una base útil y que el entrenamiento necesite muchos menos datos y muchas menos iteraciones para producir una mejora observable.

En el contexto de este TFG, esta idea es esencial. Los modelos utilizados, como \texttt{Qwen2.5-1.5B-Instruct}, \texttt{Qwen2.5-7B-Instruct} y \texttt{Qwen2.5-14B-Instruct}, ya han sido preentrenados y adaptados previamente para seguir instrucciones. El trabajo no consiste en construir un modelo desde cero, sino en estudiar cómo se puede adaptar un modelo ya existente, compararlo con versiones de distinto tamaño y medir si esa adaptación compensa desde el punto de vista de calidad y coste de inferencia.

\subsubsection{Qué ocurre durante un paso de fine tuning}

Un paso de \textit{fine tuning} puede entenderse como una repetición controlada del mismo ciclo básico de entrenamiento neuronal. Primero, se prepara una entrada $x$ y una respuesta objetivo $y$. Después, el modelo genera una distribución de probabilidad sobre los tokens de salida. A continuación, se calcula una pérdida, normalmente entropía cruzada token a token, que mide la distancia entre lo que el modelo ha predicho y lo que debería haber generado. Finalmente, se ejecuta la retropropagación y el optimizador actualiza los parámetros entrenables.

La Figura~\ref{fig:fine_tuning_ciclo} resume este proceso de forma conceptual. Lo importante no es el detalle de una implementación concreta, sino la estructura general del ciclo. Esta misma lógica se utilizará después cuando el \textit{student} aprenda a partir de señales generadas por el \textit{teacher}.

\begin{figure}[H]
\centering
\setlength{\fboxsep}{8pt}
\fbox{
\begin{minipage}{0.92\textwidth}
\centering
\footnotesize
\begin{tabularx}{0.98\textwidth}{>{\centering\arraybackslash}X c >{\centering\arraybackslash}X c >{\centering\arraybackslash}X}
\textbf{Datos de ajuste} & $\Longrightarrow$ &
\textbf{Forward del modelo} & $\Longrightarrow$ &
\textbf{Cálculo de pérdida} \\

Prompt y respuesta objetivo &
&
Predicción token a token &
&
Comparación con la respuesta esperada
\end{tabularx}

\vspace{0.7em}

\begin{tabularx}{0.98\textwidth}{>{\centering\arraybackslash}X c >{\centering\arraybackslash}X c >{\centering\arraybackslash}X}
\textbf{Actualización} & $\Longleftarrow$ &
\textbf{Backward} & $\Longleftarrow$ &
\textbf{Gradientes} \\

AdamW modifica los parámetros entrenables &
&
Retropropagación por la red &
&
Derivadas de la pérdida respecto a los pesos
\end{tabularx}
\end{minipage}
}
\caption{Ciclo conceptual de un paso de \textit{fine tuning}. El modelo recibe ejemplos de ajuste, calcula predicciones, evalúa una pérdida y actualiza los parámetros entrenables mediante retropropagación. Fuente: elaboración propia.}
\label{fig:fine_tuning_ciclo}
\end{figure}

Este ciclo es común a muchas variantes de ajuste. Lo que cambia entre unas y otras no es la lógica general, sino qué datos se usan, qué pérdida se optimiza y qué parámetros se permiten actualizar. Esta observación permite conectar el \textit{fine tuning} con la destilación. Cuando un \textit{student} se entrena con respuestas generadas por un \textit{teacher}, el procedimiento técnico sigue siendo un ajuste supervisado del modelo pequeño, aunque la fuente de las respuestas no sea humana sino sintética.

\subsubsection{Regímenes habituales de adaptación}

No todo \textit{fine tuning} persigue el mismo objetivo. En la literatura y en la práctica suelen distinguirse varios regímenes, porque cada uno modifica el modelo de una manera distinta.

El primero es el \textbf{continual pretraining}. En este caso se mantiene la tarea general de predecir el siguiente token, pero se continúa entrenando el modelo sobre un corpus nuevo, normalmente especializado en un dominio concreto. Sirve para reforzar conocimiento o estilo, pero no necesariamente para enseñar al modelo a seguir instrucciones con un formato específico.

El segundo es el \textbf{supervised fine tuning}. Aquí el modelo se entrena con pares explícitos de entrada y salida esperada. Es el régimen natural para modelos de instrucciones, diálogo y tareas donde existe una respuesta objetivo. En este TFG es especialmente relevante porque muchas formas de destilación pueden verse como un \textit{supervised fine tuning} del \textit{student} sobre datos generados por el \textit{teacher}.

El tercero es el \textbf{preference tuning}. En este caso el modelo no aprende solo a imitar una respuesta concreta, sino a preferir unas respuestas frente a otras. Métodos como RLHF o DPO pertenecen a esta familia. Aunque este TFG no entrena directamente con preferencias, es útil tenerlo presente porque muchos modelos \textit{Instruct} modernos ya incorporan etapas de este tipo antes de ser publicados.

La Tabla~\ref{tab:regimenes-finetuning} resume estas tres formas de adaptación.

\begin{table}[H]
\centering
\caption{Regímenes habituales de adaptación de LLMs. Fuente: elaboración propia.}
\label{tab:regimenes-finetuning}
\footnotesize
\begin{tabularx}{\textwidth}{>{\raggedright\arraybackslash}p{3.3cm} >{\raggedright\arraybackslash}X >{\raggedright\arraybackslash}p{4.0cm}}
\toprule
\textbf{Régimen} & \textbf{Idea principal} & \textbf{Relación con el TFG} \\
\midrule
Continual pretraining & Continuar la predicción del siguiente token sobre un corpus especializado & Se menciona como contexto, pero no es el foco del trabajo \\
\addlinespace
Supervised fine tuning & Entrenar con pares de entrada y respuesta esperada & Es el esquema más cercano a varias recetas de destilación del \textit{student} \\
\addlinespace
Preference tuning & Aprender a preferir unas respuestas frente a otras & No se entrena directamente, aunque puede estar presente en los modelos \textit{Instruct} de partida \\
\bottomrule
\end{tabularx}
\end{table}

\subsubsection{Full fine tuning}

En el \textit{full fine tuning}, todos los parámetros del modelo quedan en modo entrenable. Esto incluye embeddings, matrices de atención, bloques \textit{feed forward}, normalizaciones y cabecera de salida. A diferencia de los métodos eficientes como LoRA, aquí no se añade una pequeña corrección externa, sino que el modelo completo puede modificar internamente sus representaciones.

Esta es la opción más expresiva. Si el modelo necesita ajustar patrones profundos de razonamiento, formato, estilo o uso de conocimiento, el \textit{full fine tuning} ofrece más libertad que una adaptación de bajo rango. También es la opción que más se aproxima a decir que el modelo pequeño ha sido realmente reentrenado sobre el nuevo objetivo, aunque partiendo de una base ya preentrenada.

La contrapartida es el coste. Para entrenar todos los parámetros no basta con almacenar los pesos del modelo. También hay que almacenar gradientes, estados del optimizador y activaciones intermedias necesarias para el \textit{backward}. Una aproximación útil de la memoria de entrenamiento es
\[
M_{\mathrm{train}}
\approx
M_{\mathrm{pesos}}
+
M_{\mathrm{gradientes}}
+
M_{\mathrm{optimizador}}
+
M_{\mathrm{activaciones}}.
\]
En optimizadores como AdamW, los estados del optimizador pueden ocupar más memoria que los propios pesos, porque se mantienen momentos adicionales por cada parámetro entrenable. Por eso, un modelo que cabe en inferencia en una GPU puede no caber en entrenamiento completo en esa misma GPU.

La Tabla~\ref{tab:memoria-full-ft} resume las principales partidas de memoria durante un \textit{full fine tuning}.

\begin{table}[H]
\centering
\caption{Principales componentes de memoria durante un \textit{full fine tuning}. Fuente: elaboración propia.}
\label{tab:memoria-full-ft}
\footnotesize
\begin{tabularx}{\textwidth}{>{\raggedright\arraybackslash}p{3.0cm} >{\raggedright\arraybackslash}X >{\raggedright\arraybackslash}p{3.8cm}}
\toprule
\textbf{Componente} & \textbf{Función} & \textbf{Impacto práctico} \\
\midrule
Pesos del modelo & Parámetros actuales que definen el comportamiento de la red & Crecen linealmente con el número de parámetros \\
\addlinespace
Gradientes & Derivadas de la pérdida respecto a cada parámetro entrenable & En \textit{full fine tuning} existen para todo el modelo \\
\addlinespace
Estados del optimizador & Momentos internos usados por AdamW u otros optimizadores & Suelen ser una de las mayores fuentes de consumo de memoria \\
\addlinespace
Activaciones & Valores intermedios guardados durante el forward para calcular el backward & Crecen con batch size, longitud de secuencia y profundidad del modelo \\
\bottomrule
\end{tabularx}
\end{table}

Este coste explica por qué el \textit{full fine tuning} es más difícil de ejecutar que LoRA. Requiere más memoria, más tiempo de GPU y una configuración experimental más cuidadosa. También reduce la cantidad de pruebas que pueden realizarse dentro de una cuota de cómputo limitada. En un entorno como MareNostrum 5, esto se traduce directamente en más GPU hora, más tiempo de cola y más necesidad de cerrar bien la configuración antes de lanzar un entrenamiento largo.

Aun así, para este TFG el \textit{full fine tuning} tiene un papel metodológico muy importante. La línea final de destilación se apoya precisamente en un ajuste completo del \textit{student}, porque se quería evaluar una configuración fuerte, con la máxima capacidad de adaptación razonable dentro de los recursos disponibles. En otras palabras, LoRA sirve como herramienta eficiente para iterar y explorar, pero el \textit{full fine tuning} permite probar hasta dónde puede llegar realmente el modelo pequeño cuando todos sus parámetros pueden adaptarse a las señales del \textit{teacher}.

Esta distinción será importante al interpretar los resultados. Si una rama basada en LoRA no mejora lo suficiente, puede deberse a que la señal de destilación no es adecuada, pero también a que la adaptación de bajo rango no tiene capacidad suficiente para absorberla. En cambio, si una rama de \textit{full fine tuning} tampoco mejora, la conclusión es más fuerte, porque se ha probado una configuración con mayor libertad de actualización.

\subsubsection{Riesgos del full fine tuning}

El \textit{full fine tuning} no es solo más caro, también es más delicado. Al actualizar todos los parámetros, el modelo puede adaptarse con mucha fuerza al nuevo corpus, pero esa misma flexibilidad introduce riesgos.

El primero es el \textbf{sobreajuste}. Si el corpus de ajuste es pequeño o poco diverso, el modelo puede aprender patrones demasiado específicos del conjunto de entrenamiento y no generalizar bien al conjunto de evaluación. En este TFG esto es especialmente relevante en benchmarks como GSM8K y MATH, donde una mejora real debe observarse en el conjunto de test y no solo en la pérdida de entrenamiento.

El segundo riesgo es el \textbf{olvido catastrófico}. Al modificar todos los pesos, el modelo puede degradar capacidades generales que no aparecen representadas en el corpus de ajuste. Si el ajuste se centra únicamente en razonamiento matemático, por ejemplo, podría mejorar en ese dominio y empeorar en otros usos generales. Este riesgo no invalida el experimento, pero obliga a interpretar bien el alcance de la mejora obtenida.

El tercer riesgo es la \textbf{sensibilidad a hiperparámetros}. Una tasa de aprendizaje demasiado alta puede deteriorar rápidamente un modelo ya preentrenado. Un número excesivo de épocas puede reforzar el sobreajuste. Un formato inconsistente de datos puede enseñar patrones de salida no deseados. Por eso, cuando se emplea \textit{full fine tuning}, la calidad del dataset, el formato de las respuestas y la estabilidad de la configuración experimental son tan importantes como el propio tamaño del modelo.

\subsubsection{PEFT y la motivación de LoRA}

Las técnicas de \textit{parameter efficient fine tuning} parten de una idea empírica sencilla. En muchos casos, la diferencia entre un modelo base y su versión adaptada no necesita expresarse como una modificación densa de todos los parámetros. Basta con aprender una corrección relativamente pequeña, concentrada en un subespacio de baja dimensión. Esta observación ha dado lugar a una familia de métodos que congelan la mayor parte del modelo y entrenan solo una fracción reducida de parámetros.

Dentro de esta familia se encuentran varias propuestas, pero la más influyente y la más relevante para este TFG es \textbf{LoRA}. Su interés radica en que reduce drásticamente el coste de entrenamiento sin introducir una penalización significativa en inferencia una vez que la adaptación se fusiona con el modelo base.

\subsubsection{LoRA como adaptación de bajo rango}
\label{sec:finetuning-lora}

La idea central de LoRA es aproximar la actualización de una matriz de pesos mediante una corrección de bajo rango~\cite{hu2021lora}. Considérese una capa lineal con peso
\[
W_0 \in \mathbb{R}^{d_{\mathrm{out}}\times d_{\mathrm{in}}}.
\]
En el \textit{full fine tuning}, esta matriz se actualizaría libremente. En LoRA, en cambio, se mantiene congelada y se aprende una corrección
\[
\Delta W \approx BA,
\]
donde
\[
A \in \mathbb{R}^{r\times d_{\mathrm{in}}},
\qquad
B \in \mathbb{R}^{d_{\mathrm{out}}\times r},
\qquad
r \ll \min(d_{\mathrm{in}},d_{\mathrm{out}}).
\]
De este modo, la transformación lineal pasa a escribirse como
\[
y
=
W_0 x
+
\frac{\alpha}{r}BAx,
\]
donde $r$ es el rango de la adaptación y $\alpha$ es un factor de escala que controla la magnitud efectiva de la corrección.

La ventaja es inmediata. En lugar de entrenar $d_{\mathrm{out}}d_{\mathrm{in}}$ parámetros para esa matriz, LoRA solo añade
\[
r(d_{\mathrm{in}} + d_{\mathrm{out}})
\]
parámetros entrenables. Cuando $r$ es pequeño, la reducción es de varios órdenes de magnitud. Esto se traduce en menos gradientes, menos estados del optimizador y un consumo de memoria mucho menor durante el entrenamiento.

La Figura~\ref{fig:lora-correccion} representa la diferencia conceptual entre \textit{full fine tuning} y LoRA.

\begin{figure}[H]
\centering
\setlength{\fboxsep}{8pt}
\fbox{%
\begin{minipage}{0.92\textwidth}
\centering
\footnotesize
% No usar display math \[...\] dentro de tabularx: el alineador & rompe el modo matemático.
\begin{minipage}[t]{0.44\textwidth}\centering
\textbf{Full fine tuning}\\[0.4em]
Todos los pesos del modelo pueden cambiar
\[
W_0 \rightarrow W_0 + \Delta W
\]
\[
\Delta W\ \text{es libre}
\]
\end{minipage}\hfill
\begin{minipage}[t]{0.44\textwidth}\centering
\textbf{LoRA}\\[0.4em]
El peso original se congela y se añade una corrección pequeña
\[
W_0 x + \frac{\alpha}{r}BAx
\]
\[
\Delta W \approx BA,\quad r \ll d
\]
\end{minipage}
\end{minipage}%
}%
\caption{Comparación conceptual entre \textit{full fine tuning} y LoRA. En el primer caso todos los pesos pueden actualizarse. En LoRA, el peso original se mantiene congelado y se aprende una corrección de bajo rango. Fuente: elaboración propia a partir de~\cite{hu2021lora}.}
\label{fig:lora-correccion}
\end{figure}

Además, los pesos originales del modelo permanecen intactos. Esto tiene dos consecuencias relevantes. La primera es que el riesgo de alterar fuertemente capacidades generales del modelo base suele ser menor que en un \textit{full fine tuning} agresivo. La segunda es que, una vez entrenadas las matrices de LoRA, la corrección puede fusionarse con el peso original
\[
W_{\mathrm{merged}} = W_0 + \frac{\alpha}{r}BA,
\]
de manera que en inferencia el modelo resultante se comporta como un modelo convencional, sin necesidad de ejecutar una ruta especial de \textit{forward}.

La Tabla~\ref{tab:ft-vs-lora} resume las diferencias más importantes entre \textit{full fine tuning} y LoRA.

\begin{table}[H]
\centering
\caption{Comparación conceptual entre \textit{full fine tuning} y LoRA. Fuente: elaboración propia a partir de~\cite{hu2021lora}.}
\label{tab:ft-vs-lora}
\footnotesize
\begin{tabularx}{\textwidth}{>{\raggedright\arraybackslash}p{3.0cm} >{\raggedright\arraybackslash}X >{\raggedright\arraybackslash}X}
\toprule
\textbf{Aspecto} & \textbf{Full fine tuning} & \textbf{LoRA} \\
\midrule
Parámetros entrenables & Todos los parámetros del modelo & Solo las matrices de bajo rango añadidas a ciertas capas \\
\addlinespace
Coste de memoria en entrenamiento & Muy alto, por gradientes y estados del optimizador sobre todo el modelo & Mucho menor, porque el conjunto entrenable es reducido \\
\addlinespace
Capacidad de adaptación & Máxima expresividad dentro del modelo elegido & Menor expresividad, aunque suele ser suficiente en adaptaciones moderadas \\
\addlinespace
Complejidad experimental & Mayor, especialmente en modelos grandes & Menor, permite iterar más rápido y con menos recursos \\
\addlinespace
Impacto en inferencia & Ninguno una vez entrenado & Ninguno una vez fusionado con el modelo base \\
\addlinespace
Papel en el TFG & Rama fuerte para maximizar capacidad de adaptación del \textit{student} & Rama eficiente para explorar configuraciones con bajo coste \\
\bottomrule
\end{tabularx}
\end{table}

\subsubsection{Relación con la destilación y con el TFG}

La conexión entre esta sección y la siguiente es directa. Desde un punto de vista práctico, muchas recetas de \textit{knowledge distillation} pueden implementarse como una forma de \textit{fine tuning} del \textit{student} sobre señales procedentes del \textit{teacher}. Por eso resulta imprescindible entender antes qué significa ajustar un modelo pequeño, cuánto cuesta hacerlo y qué alternativas existen para reducir ese coste.

En el TFG aparecen dos formas de adaptación especialmente relevantes. Por un lado, ramas de ajuste eficientes basadas en LoRA, útiles para iterar rápidamente y explorar configuraciones de bajo coste. Por otro, una rama de \textit{full fine tuning}, concebida como referencia de máxima capacidad adaptativa dentro de los recursos disponibles. La comparación entre ambas no solo tiene interés práctico, sino también metodológico, porque muestra con claridad el compromiso entre expresividad, coste de entrenamiento y viabilidad experimental.

Con este marco ya establecido, la destilación de conocimiento puede abordarse en el apartado siguiente no como una técnica aislada, sino como una pieza concreta dentro de un ecosistema más amplio de métodos para mejorar la eficiencia de los modelos y de su despliegue.

\section{Knowledge Distillation}
\label{sec:knowledge-distillation}

La \textit{knowledge distillation} (KD) es la técnica que hace posible la cascada de este TFG, porque sin un student que imite razonablemente bien al teacher, la idea de escalar al grande solo cuando haga falta pierde todo su sentido. En esta sección se explica qué es la destilación, cómo se formula matemáticamente y qué variantes se usan hoy en día, con especial atención a la que se emplea en los experimentos.

\subsection{Definición y motivación}

La idea central fue propuesta por Hinton, Vinyals y Dean en 2015~\cite{hinton2015distilling} para modelos de clasificación, y ha sido extendida desde entonces a prácticamente cualquier tipo de red neuronal, incluidos los LLMs generativos. El punto de partida es casi filosófico. Cuando un modelo grande ha sido entrenado sobre un corpus masivo, la información que ha adquirido no reside únicamente en cuál es su predicción correcta, sino también (y esto es lo importante) en \emph{cómo distribuye la probabilidad entre las alternativas incorrectas}. Esa distribución es mucho más rica que una simple etiqueta \emph{one-hot}.

Para entender por qué, conviene detenerse un momento en un ejemplo. Supongamos que un modelo grande se enfrenta a un problema de aritmética que debería resolverse con 42. Una etiqueta dura solo diría ``la respuesta es 42''. El teacher, en cambio, dirá algo más parecido a: ``mi respuesta es 42 con probabilidad 0,85; 24 con 0,08; 12 con 0,04; el resto del vocabulario es esencialmente cero''. Esa estructura secundaria no es ruido: contiene información sobre qué errores son \emph{plausibles} (24 y 12 son números del orden correcto y con quizá algún paso intermedio mal hecho) y cuáles son \emph{imposibles} (99 o una letra). Hinton llamó a esa información los \textit{soft labels} o \textit{dark knowledge}, por contraste con las hard labels del entrenamiento clásico. Un modelo pequeño que aprenda a imitar esas distribuciones suaves está recibiendo, en cada ejemplo, una señal mucho más densa que la que le daría una etiqueta one-hot. En lugar de un único bit de información por token, recibe una descripción probabilística completa sobre el vocabulario entero.

La consecuencia empírica es significativa. Un student entrenado con distilación suele recuperar una fracción muy apreciable de la calidad del teacher, por encima de la que se conseguiría entrenando el mismo modelo pequeño desde cero con los mismos datos. Esto se explica por dos factores complementarios. Primero, el student accede a una \emph{función objetivo mejor condicionada}: la distribución del teacher ya es una aproximación razonable del comportamiento deseado, y encajar en ella es más fácil que deducirla solo a partir de etiquetas binarias. Segundo, las soft labels contienen implícitamente conocimiento sobre la \emph{estructura} del problema: qué clases se parecen entre sí, qué errores son ``cercanos'' y cuáles ``lejanos''. El student no tiene que reconstruir esa geometría por sí solo, la recibe ya extraída.

Esta propiedad es especialmente valiosa en LLMs, donde entrenar un modelo de decenas de miles de millones de parámetros desde cero con tokens suficientes cuesta decenas o centenares de miles de horas de GPU. Destilar un student a partir de un teacher ya entrenado, sobre un corpus relativamente modesto, es órdenes de magnitud más barato y, si se hace con cuidado, conserva buena parte del comportamiento del modelo grande, especialmente en dominios bien cubiertos por los datos de destilación. Por eso la destilación aparece una y otra vez en la literatura reciente como pieza central del pipeline de entrega de modelos a producción, y por eso es la técnica habilitadora de la cascada en este TFG.

\subsection{Formulación clásica de la destilación por respuesta}

La versión original de Hinton se formula sobre las distribuciones de salida del teacher y del student. Sea un input $x$ y su etiqueta verdadera $y$. Sea $z^t(x)$ el vector de logits del teacher y $z^s(x)$ el del student. Sobre esos logits se aplica un softmax con \emph{temperatura} $T$:
\[
p^t_i(x; T) = \frac{\exp\!\big(z^t_i(x)/T\big)}{\sum_j \exp\!\big(z^t_j(x)/T\big)}, \qquad
p^s_i(x; T) = \frac{\exp\!\big(z^s_i(x)/T\big)}{\sum_j \exp\!\big(z^s_j(x)/T\big)}.
\]
La pérdida de destilación es una combinación convexa de dos términos: una pérdida de entropía cruzada contra la etiqueta real y una divergencia de Kullback--Leibler entre las distribuciones del teacher y del student,
\[
\mathcal{L}_{\mathrm{KD}}(x, y) \;=\; \alpha \, \underbrace{\mathrm{CE}\!\big(y,\, p^s(x;1)\big)}_{\text{pérdida supervisada}} \;+\; (1-\alpha)\, T^2 \, \underbrace{\mathrm{KL}\!\big(p^t(x;T)\,\|\,p^s(x;T)\big)}_{\text{pérdida de destilación}}.
\]

La expresión parece inocente pero cada pieza tiene un motivo, y conviene desgranarla con cuidado porque los mismos ingredientes (KL, temperatura, balance) aparecen modificados en todas las variantes modernas.

\paragraph{Qué mide la divergencia KL y por qué aquí.} La divergencia de Kullback--Leibler entre dos distribuciones $p$ y $q$ sobre el mismo soporte se define como
\[
\mathrm{KL}(p\,\|\,q) \;=\; \sum_{i} p_i \log\!\frac{p_i}{q_i}.
\]
No es simétrica ($\mathrm{KL}(p\|q)\neq \mathrm{KL}(q\|p)$) y admite una interpretación muy natural: mide cuántos bits (o \emph{nats}) de información se pierden de media si se usa $q$ para codificar muestras que en realidad provienen de $p$. Al ponerla en la pérdida como $\mathrm{KL}(p^t\|p^s)$, estamos imponiendo al student que su distribución $p^s$ se acerque a la del teacher $p^t$ en el sentido de \emph{cubrir la masa de probabilidad del teacher}: si el teacher asigna probabilidad alta a una clase, el student será penalizado fuertemente por no seguirlo; si el teacher asigna probabilidad cero, el término correspondiente desaparece y no impone restricción alguna. Esta asimetría es deseable en distilación, porque lo que nos importa es que el student \emph{no olvide} lo que el teacher cree probable, no al revés.

\paragraph{El papel de la temperatura $T$.} Al dividir los logits por $T>1$ antes del softmax, la distribución se ``suaviza'': las probabilidades más altas bajan y las más bajas suben. Es un mecanismo crucial porque en modelos bien entrenados la distribución del teacher suele estar muy concentrada en la clase correcta (por ejemplo, 0,99 en una clase y 0,01 repartido en el resto). En esa situación, el softmax original \emph{aplasta} toda la información secundaria: las probabilidades relativas entre las clases incorrectas son indistinguibles numéricamente y no aportan señal apenas. Aumentar $T$ tiene el efecto de magnificar esas diferencias y ponerlas en un rango donde el student puede realmente aprenderlas. A cambio, al suavizar hacemos también que los gradientes del student se vuelvan más pequeños (las diferencias entre probabilidades son más sutiles), y por eso el segundo término se multiplica por $T^2$, para compensar esta reducción y mantener el orden de magnitud del gradiente invariante con $T$~\cite{hinton2015distilling}. En la práctica, valores razonables son $T\in[2,10]$ según la tarea y el modelo; $T$ muy grande tiende a ``aplastarlo todo'' por exceso y deja de aportar información.

\paragraph{El papel de $\alpha$.} Es el hiperparámetro que balancea cuánto pesa aprender de la verdad objetiva (las hard labels) frente a cuánto pesa imitar al teacher (las soft labels). Valores bajos de $\alpha$ significan que el student casi solo imita al teacher, lo que le permite heredar comportamiento incluso en ejemplos donde la etiqueta dura no es informativa. Valores altos lo anclan a la verdad, útil cuando el teacher está sesgado o se equivoca sistemáticamente en ciertos casos. Un valor típico en la literatura es $\alpha \in [0{,}1, 0{,}5]$, pero depende fuertemente de la tarea y de la calidad del teacher.

\paragraph{Qué está aprendiendo realmente el student.} Puede resultar útil pensar la pérdida KD como un objetivo con dos ``supervisores'' complementarios. Las hard labels le dicen al student \emph{qué} responder; las soft labels del teacher le enseñan, además, \emph{cómo se parece cada posible respuesta a la correcta}. Esa segunda señal es estructural: fija cómo debe organizar el student su espacio de salida, qué errores debe considerar ``cercanos'' a la respuesta correcta y cuáles ``lejanos''. Empíricamente, esta estructura es transferible a ejemplos nuevos (generaliza mejor) y es una de las razones de que KD sea más que un simple \textit{label smoothing}. De hecho, cuando el teacher es de muy alta calidad, distilar con $\alpha$ cercano a cero y $T$ elevado puede superar en generalización al entrenamiento supervisado clásico incluso cuando el student tiene la misma capacidad que el teacher.

\begin{figure}[H]
\centering
\resizebox{\textwidth}{!}{%
\begin{tikzpicture}[
>=Latex,
mdl/.style={draw, rounded corners=2pt, align=center, minimum width=2.6cm, minimum height=9mm, font=\footnotesize, thick},
logit/.style={draw, rounded corners=2pt, align=center, minimum width=2.1cm, minimum height=7mm, font=\footnotesize, fill=gray!10},
loss/.style={draw, rounded corners=2pt, align=center, minimum width=3.2cm, minimum height=8mm, font=\footnotesize, thick},
arr/.style={->, thick}
]
% Entrada
\node[font=\footnotesize] (x1) at (0,0) {$x$};

% Modelos
\node[mdl, fill=bcteacher!40] (teacher) at (2.3, 1.3) {Teacher\\(gran modelo)};
\node[mdl, fill=bcstudent!40] (student) at (2.3, -1.3) {Student\\(modelo pequeño)};

% Logits / probabilidades
\node[logit] (pt) at (5.2, 1.3) {$p^t(x;T)$};
\node[logit] (ps) at (5.2, -1.3) {$p^s(x;T)$};

% Etiqueta dura
\node[logit, fill=bcdata!40] (y) at (5.2, -3.0) {etiqueta $y$};

% Losses parciales
\node[loss, fill=bcrouter!30] (kl) at (9.0, 1.3) {$T^{2}\,\mathrm{KL}(p^{t}\|p^{s})$};
\node[loss, fill=bcdata!25] (ce) at (9.0, -2.3) {$\mathrm{CE}(y,\, p^{s}(x;1))$};

% Loss total
\node[loss, fill=bcmeasure!35, minimum width=4.4cm] (sum) at (13.3, -0.5) {$\mathcal{L}=\alpha\,\mathrm{CE}+(1{-}\alpha)\,T^{2}\,\mathrm{KL}$};

% Flechas de entrada a modelos
\draw[arr] (x1) -- (teacher.west);
\draw[arr] (x1) -- (student.west);

% Modelos a logits
\draw[arr] (teacher.east) -- (pt.west);
\draw[arr] (student.east) -- (ps.west);

% pt -> KL (recta)
\draw[arr] (pt.east) -- (kl.west);
% ps -> KL (sube, pasando a la derecha de pt)
\draw[arr] (ps.east) -- ++(0.6,0) |- (kl.south);

% y -> CE
\draw[arr] (y.east) -- (ce.west);
% ps -> CE (pasa por debajo, dashed indica T=1)
\draw[arr, dashed] (ps.south) -- ++(0,-0.55) -| (ce.north);

% KL y CE -> L
\draw[arr] (kl.east) -| (sum.north);
\draw[arr] (ce.east) -| (sum.south);
\end{tikzpicture}%
}%
\caption{Esquema de la destilación clásica por respuesta: el teacher y el student procesan la misma entrada, y la pérdida del student combina una CE contra la etiqueta dura con una KL contra la distribución suavizada del teacher. La línea discontinua indica que el término CE usa el student evaluado a temperatura 1. Fuente: elaboración propia.}
\label{fig:kd-esquema}
\end{figure}

\subsection{De los modelos clasificadores a los modelos generativos}

La formulación anterior se pensó originalmente para clasificadores: una entrada, una salida categórica. En LLMs generativos la situación es distinta porque la ``salida'' no es una etiqueta única, sino una secuencia de tokens $y = (y_1, y_2, \dots, y_T)$ generada de forma autoregresiva. Esto obliga a repensar qué significa imitar al teacher.

Una opción es aplicar la KD clásica \emph{token a token}: en cada posición $i$ de una secuencia fijada, el teacher produce una distribución $p^t(\cdot \mid y_{<i}, x)$ sobre el siguiente token y el student se entrena para aproximar esa distribución con una KL. Esta variante se conoce como \textbf{word-level KD} o \textbf{on-logits KD} y es la extensión más directa de Hinton al entorno autoregresivo. Tiene la ventaja de transferir todo el detalle probabilístico del teacher, pero presenta inconvenientes prácticos en nuestro contexto: exige mantener el teacher \emph{cargado durante todo el entrenamiento} del student, lo que implica que dos modelos muy grandes deben convivir en memoria de GPU; requiere, además, que el vocabulario del student sea idéntico al del teacher (o al menos compatible mediante una proyección), porque las distribuciones se comparan token a token sobre el mismo vocabulario.

La alternativa, descrita por Kim y Rush~\cite{kim2016sequence}, es la \textbf{destilación a nivel de secuencia} o \textit{sequence-level KD}. En ella, el teacher primero genera respuestas completas $\hat{y}^t(x)$ para cada entrada $x$ del corpus y, después, el student se entrena para \emph{producir} esas mismas respuestas mediante entropía cruzada estándar, sin acceso en línea a las distribuciones del teacher. La pérdida pasa a ser simplemente
\[
\mathcal{L}_{\mathrm{seq-KD}}(x) \;=\; -\sum_{i=1}^{T} \log p^s\!\big(\hat{y}^t_i(x)\,\big|\,\hat{y}^t_{<i}(x),\, x\big),
\]
es decir, un fine-tuning supervisado clásico sobre un corpus sintético generado por el teacher. Conceptualmente, se puede ver como una aproximación de la KD token a token en la que la distribución del teacher se sustituye por una \emph{muestra} suya (la respuesta generada). Se pierde la riqueza de las soft labels (ahora cada posición vuelve a ser una etiqueta \emph{hard} marcada por el teacher), pero se ganan propiedades de ingeniería muy valiosas para trabajar con LLMs grandes.

En primer lugar, \textbf{el teacher se usa una única vez} para generar el corpus sintético y se descarga. El entrenamiento posterior del student es un SFT estándar que cabe en una sola GPU y reutiliza todo el ecosistema habitual de fine-tuning (transformers, TRL, PEFT). En segundo lugar, \textbf{no se requiere compatibilidad de vocabulario ni de arquitectura} entre teacher y student: basta con que ambos produzcan texto. Esto abre la puerta a destilar entre familias de modelos distintas o incluso a destilar salidas de modelos cerrados vía API. En tercer lugar, \textbf{el corpus sintético es reutilizable}: distintos students, distintas configuraciones de hiperparámetros o distintas mezclas de dominios se pueden entrenar sobre el mismo corpus, lo que abarata considerablemente las iteraciones experimentales.

El coste se paga en otro lado: el student aprende a reproducir las respuestas del teacher, no su \emph{distribución}. Si una pregunta tiene varias soluciones válidas, el student solo ``ve'' la que tocó generar el teacher y pierde información sobre alternativas razonables. Empíricamente, la literatura reciente sobre LLMs muestra que esta pérdida es a menudo aceptable cuando el corpus sintético es suficientemente grande y diverso, y que los beneficios prácticos compensan ampliamente el coste. Por ese motivo, la destilación a nivel de secuencia es hoy el esquema dominante en sistemas reales, y es también el que se adopta en este TFG.

\begin{figure}[H]
\centering
\resizebox{\textwidth}{!}{%
\begin{tikzpicture}[
>=Latex,
mdl/.style={draw, rounded corners=3pt, align=center, minimum width=30mm, minimum height=11mm, font=\footnotesize, thick},
data/.style={draw, rounded corners=2pt, align=center, minimum width=28mm, minimum height=9mm, font=\footnotesize, fill=bcdata!30},
arr/.style={->, thick}
]
% Fase 1: Generación del corpus
\node[font=\bfseries\footnotesize] at (0, 2.6) {Fase 1: Generación};
\node[data] (prompts) at (0, 1.5) {Prompts\\(GSM8K, MATH)};
\node[mdl, fill=bcteacher!40] (teacher) at (0, 0) {Teacher\\(7B / 14B)};
\node[data, fill=bcmeasure!25] (corpus) at (0, -1.5) {Corpus sintético\\$\{(x_i, \hat{y}^t_i)\}$};
\draw[arr] (prompts) -- (teacher);
\draw[arr] (teacher) -- (corpus);

% Fase 2: SFT del student
\node[font=\bfseries\footnotesize] at (6.5, 2.6) {Fase 2: SFT del student};
\node[mdl, fill=bcstudent!40] (student) at (6.5, 0) {Student\\(1.5B)};
\node[data, fill=bcrouter!25] (trained) at (6.5, -1.5) {Student\\destilado};

\draw[arr] (corpus.east) -- ++(1.2,0) |- (student.west) node[pos=0.25, above, font=\scriptsize] {entrenamiento SFT};
\draw[arr] (student) -- (trained);

% Anotaciones
\node[font=\scriptsize, align=center, text=gray] at (0, -2.5) {El teacher se descarga\\tras esta fase};
\node[font=\scriptsize, align=center, text=gray] at (6.5, -2.5) {Reutiliza tooling estándar\\(transformers, TRL, PEFT)};
\end{tikzpicture}%
}%
\caption{Pipeline de destilación \textit{sequence-level} adoptado en este TFG. En la primera fase, el teacher genera respuestas completas sobre los prompts de los benchmarks. En la segunda, el student se ajusta vía SFT sobre ese corpus sintético. El teacher no interviene en la segunda fase, lo que permite iterar sobre el student con una sola GPU. Fuente: elaboración propia.}
\label{fig:seq-kd-pipeline}
\end{figure}

\subsection{Otras variantes relevantes}

La taxonomía completa es mucho más rica. La Tabla~\ref{tab:kd-variantes} recoge las variantes que es útil conocer para situar el trabajo, agrupadas por la naturaleza de la señal que se transfiere.

\begin{table}[H]
\centering
\caption{Principales variantes de destilación de conocimiento en modelos neuronales. Fuente: elaboración propia a partir de~\cite{hinton2015distilling, sanh2019distilbert, gou2021survey, agarwal2024gkd}.}
\label{tab:kd-variantes}
\footnotesize
\begin{tabularx}{\textwidth}{>{\raggedright\arraybackslash}p{3cm} >{\raggedright\arraybackslash}X >{\raggedright\arraybackslash}p{3.5cm}}
\toprule
\textbf{Variante} & \textbf{Qué se transfiere} & \textbf{Ejemplo representativo} \\
\midrule
Response-based (\textit{word-level KD}) & Distribución del teacher sobre el vocabulario, token a token & Hinton et al.~\cite{hinton2015distilling}; DistilBERT~\cite{sanh2019distilbert} \\
\addlinespace
Sequence-level KD & Secuencia completa generada por el teacher, usada como etiqueta para un SFT del student & Kim \& Rush~\cite{kim2016sequence}; Alpaca/Vicuna~\cite{taori2023alpaca} \\
\addlinespace
Feature-based & Activaciones intermedias (representaciones ocultas) en capas seleccionadas & FitNets, hidden-state distillation \\
\addlinespace
Relation-based & Relaciones entre ejemplos (distancias o similitudes entre representaciones) & Relational KD, RKD \\
\addlinespace
On-policy / generative KD & El student genera, el teacher puntúa, y la pérdida se aplica sobre las secuencias del propio student & GKD~\cite{agarwal2024gkd} \\
\bottomrule
\end{tabularx}
\end{table}

\textbf{Response-based (word-level).} Es la destilación clásica adaptada al entorno autoregresivo: el student aprende a igualar la distribución del teacher en cada posición. Muy completa en términos de información transferida, pero cara de orquestar cuando teacher y student son ambos grandes, y restringida a pares con vocabulario compatible.

\textbf{Sequence-level.} Es la variante adoptada en este TFG y su presencia en la tabla por derecho propio refleja su importancia práctica en la era de los LLMs. El teacher genera respuestas completas y el student se ajusta vía SFT sobre ellas. Es el esquema que está detrás de la mayoría de modelos derivados que han aparecido en 2023--2024, donde modelos abiertos de unos pocos miles de millones de parámetros se entrenan sobre salidas generadas por modelos más grandes (propietarios o abiertos). En realidad, muchos de los modelos que hoy conocemos bajo la etiqueta de ``\textit{student}'' son sequence-level KD disfrazados.

\textbf{Feature-based.} El student no solo reproduce la salida final sino también activaciones intermedias en capas seleccionadas. Gana expresividad porque transfiere representaciones internas, pero necesita una correspondencia razonable entre las arquitecturas del teacher y del student. Si el student tiene menos capas o dimensiones distintas, hay que añadir proyectores entrenables que alineen los espacios. En contextos donde teacher y student son de la misma familia y difieren solo en tamaño (por ejemplo, en Qwen2.5 de 14B y 7B), feature-based es técnicamente factible pero rara vez aporta ganancias claras por encima de la sequence-level KD.

\textbf{Relation-based.} En lugar de imitar \emph{qué} dice el teacher sobre cada ejemplo, imita \emph{cómo relaciona} el teacher distintos ejemplos entre sí (distancias, similitudes, ángulos en el espacio de representación). Es útil en \textit{representation learning} (por ejemplo, para destilar encoders) pero menos directo en LLMs generativos, donde la salida es texto.

\textbf{On-policy vs off-policy.} Esta distinción es ortogonal a las anteriores y ha ganado peso en los últimos años. En la destilación \textit{off-policy} (también llamada ``por imitación'' o de corpus fijo) el teacher genera previamente y el student se entrena sobre esas secuencias. Es sencilla, reproducible y permite reutilizar el corpus entre experimentos. Pero tiene un problema conocido como \textit{train--test mismatch}: el student, una vez entrenado, genera sus propias secuencias que pueden separarse distribucionalmente de las del teacher; en esas zonas nuevas del espacio, el student nunca fue entrenado y su comportamiento puede degradarse.

En la destilación \textit{on-policy} o \textit{generative} (GKD)~\cite{agarwal2024gkd}, el student genera sus propias secuencias durante el entrenamiento, el teacher las evalúa en tiempo real y la pérdida se aplica sobre esas secuencias generadas. Esto cierra el mismatch: el student aprende precisamente en las zonas donde acaba moviéndose en inferencia. A cambio, es mucho más caro porque requiere tanto \textit{forward} del student (para generar) como \textit{forward} del teacher (para puntuar) en bucle, y exige mantener ambos modelos en memoria simultáneamente. En presupuestos de cómputo ajustados, la estrategia habitual es priorizar el off-policy y, si sobra presupuesto, añadir una fase final de refinamiento on-policy.

\subsection{Cómo encaja la destilación en este TFG}

En este TFG, el papel de la destilación es \textbf{habilitar} la cascada. Su rol no es introducir una técnica nueva de KD, sino obtener uno o más students capaces de atender un porcentaje razonable de peticiones con calidad aceptable, de modo que las políticas B y C tengan un modelo pequeño al que merezca la pena enviar tráfico antes de escalar al grande. Dentro del mapa anterior, el esquema concreto adoptado es una destilación \textbf{sequence-level, off-policy}. Esta elección es consecuencia directa de varias restricciones prácticas del proyecto que vale la pena hacer explícitas.

La primera es el \textbf{presupuesto de GPU acotado}. El corpus de destilación se genera una única vez con el teacher sobre GSM8K y MATH, y el teacher se descarga inmediatamente después. Todos los entrenamientos posteriores del student se realizan con un único modelo en GPU, lo que permite iterar con rapidez sobre configuraciones distintas sin consumir cuota de forma desproporcionada.

La segunda es la \textbf{compatibilidad de \textit{tooling}} y la simplicidad operativa. Una vez que se dispone de un corpus de pares \emph{(prompt, respuesta del teacher)}, el entrenamiento del student es un fine-tuning supervisado estándar. Esto permite aprovechar el ecosistema maduro de herramientas de SFT (transformers, TRL, PEFT) sin necesidad de infraestructura \textit{ad hoc}, y combinar la destilación con LoRA (Sec.~\ref{sec:finetuning-lora}) en las ramas de iteración rápida y con fine-tuning completo en la rama final más ambiciosa.

La tercera es la \textbf{cobertura del dominio}. Los benchmarks objetivo son matemáticos, donde la correctitud final es binaria y el razonamiento intermedio suele ser reproducible por un student si el corpus de destilación cubre los tipos de problema con suficiente diversidad. En este régimen, la literatura reciente sugiere que la pérdida de información al pasar de \emph{word-level} a \emph{sequence-level} KD es pequeña y claramente aceptable frente al ahorro de cómputo.

La cuarta es la \textbf{reusabilidad del corpus}. El mismo corpus sintético se ha usado como base para múltiples experimentos de destilación del TFG (distintas configuraciones de hiperparámetros, LoRA frente a FT completo, filtrado por corrección de respuestas, entre otros), lo que es una propiedad exclusiva del esquema off-policy.

\paragraph{Plan de contingencia.} Si el protocolo de decisión T5.3 hubiera concluido que la destilación no era viable con la cuota disponible, se habría adoptado el plan de contingencia R2 con un modelo preentrenado de la misma familia (\texttt{Qwen2.5-1.5B-Instruct} directamente, sin ajuste adicional), aceptando una cota de calidad algo más baja como precio por no tener que entrenar. Esta decisión queda explícitamente documentada y no invalida el resto del experimento, porque las políticas de cascada se pueden evaluar contra cualquier student, destilado o no. En la práctica el proyecto ha llegado a ejecutar los experimentos de destilación y, por tanto, la configuración de referencia que se reporta en la parte experimental es la del student destilado.

%%%%%%%%%%%%%%%%%%%%%%%%%%%%%%%%%%%%%%%%%%%%%%%%%%%%%%%%%%%%%%%%%
\section{Serving, cascadas de modelos e infraestructura HPC}
\label{sec:serving-teorico}

Tener un teacher y un student entrenados no basta para que la cascada funcione en la práctica. Hace falta un motor de inferencia capaz de servir ambos modelos con buena utilización de la GPU, baja latencia de primer token y throughput alto, y con una API estable sobre la que construir el router. Este bloque enlaza tres piezas que en el índice conviene tratar unidas: el \textbf{serving} moderno con vLLM~\cite{vllm2023, vllm-docs}, las \textbf{cascadas y políticas de routing} entre modelos, y el \textbf{entorno HPC} (MareNostrum~5, SLURM) donde se ejecuta el pipeline experimental.

\subsection{Serving con vLLM (memoria, paginación y batching)}

\subsubsection{El problema de memoria y la fragmentación del KV cache}

Como se ha visto en la sección anterior, el KV cache crece linealmente con el número de tokens y ocupa buena parte de la VRAM disponible. El problema se agrava cuando se sirven varias peticiones a la vez, porque cada una tiene su propio KV cache de tamaño variable y dinámico: algunas acabarán generando 50 tokens y otras 2\,000, y no se sabe a priori cuáles.

Los sistemas de serving anteriores a vLLM reservaban el KV cache como un \textbf{bloque contiguo} por petición, dimensionado al peor caso (la longitud máxima esperada). Esto produce dos formas de desperdicio. La \textbf{fragmentación interna} aparece porque cada petición reserva mucho más espacio del que acaba usando: una petición que termina con 300 tokens pero reservó espacio para 2\,048 desperdicia el 85\,\% de esa memoria. La \textbf{fragmentación externa} ocurre porque entre peticiones quedan huecos de memoria contigua que no son suficientemente grandes para una petición nueva pero sí para varias pequeñas, y el gestor clásico no es capaz de aprovecharlos.

El resultado es que, con gestión ingenua, una GPU con 80~GB efectivamente solo puede sostener una fracción reducida de ese espacio como KV cache útil, y el motor de serving acaba limitado no por cómputo sino por memoria disponible para sostener el batch. La consecuencia directa es un throughput bajo.

\subsubsection{PagedAttention y vLLM}

vLLM introduce \textbf{PagedAttention}~\cite{vllm2023}, una idea inspirada directamente en cómo los sistemas operativos gestionan la memoria virtual. La memoria del KV cache se divide en \textbf{bloques} de tamaño fijo (típicamente 16 tokens por bloque). Cada petición mantiene una tabla lógica que mapea sus tokens a bloques físicos no necesariamente contiguos. Cuando una petición crece, el motor le asigna nuevos bloques del pool libre, vengan de donde vengan. Cuando una petición termina, sus bloques vuelven al pool y pueden ser reutilizados por cualquier otra.

Los beneficios de este diseño son inmediatos. En primer lugar, se elimina la fragmentación interna porque solo se asignan bloques cuando son necesarios y no hay ``reserva al peor caso''. En segundo lugar, desaparece la fragmentación externa porque no se requiere memoria contigua: cualquier bloque libre del pool sirve para cualquier petición. En tercer lugar, PagedAttention habilita la compartición de prefijos: cuando varias peticiones comparten prompt (por ejemplo, el mismo \textit{system prompt}), los bloques correspondientes al prefijo común pueden reutilizarse sin duplicarse, empleando un mecanismo de \textit{copy-on-write} si alguna petición diverge.

El trabajo original~\cite{vllm2023} reportó mejoras de throughput de 2--4$\times$ respecto a motores de serving anteriores en los mismos escenarios, manteniendo la latencia. A día de hoy vLLM es una de las piezas más usadas en serving de LLMs abiertos en producción y es el motor de referencia en el ecosistema del BSC para LLMs.

\begin{figure}[H]
\centering
\resizebox{0.95\textwidth}{!}{%
\begin{tikzpicture}[
>=Latex,
legacyblk/.style={draw, fill=red!25, minimum height=6mm, minimum width=6mm, inner sep=0},
freeblk/.style={draw, fill=gray!20, minimum height=6mm, minimum width=6mm, inner sep=0},
paged/.style={draw, fill=blue!30, minimum height=6mm, minimum width=6mm, inner sep=0},
lbl/.style={font=\scriptsize},
rowlbl/.style={font=\footnotesize\bfseries}
]
% Fila 1: Legacy (reserva contigua con huecos)
\node[rowlbl, anchor=east] at (-0.3, 0) {Legacy};
\foreach \x/\s in {0/legacyblk, 1/legacyblk, 2/legacyblk, 3/freeblk, 4/freeblk, 5/legacyblk, 6/legacyblk, 7/freeblk, 8/freeblk, 9/freeblk, 10/legacyblk, 11/legacyblk}{
  \node[\s] at (\x*0.72, 0) {};
}
\node[lbl, anchor=west] at (11*0.72+0.55, 0) {$\sim$40\% útil};

% Fila 2: PagedAttention (todos los bloques aprovechables)
\node[rowlbl, anchor=east] at (-0.3, -1.1) {PagedAttn};
\foreach \x in {0,1,2,3,4,5,6,7,8,9,10,11}{
  \node[paged] at (\x*0.72, -1.1) {};
}
\node[lbl, anchor=west] at (11*0.72+0.55, -1.1) {$\sim$90\% útil};

% Leyenda de colores
\node[legacyblk] at (0.5, -2.3) {};
\node[lbl, anchor=west] at (0.9, -2.3) {ocupado};
\node[freeblk] at (3.0, -2.3) {};
\node[lbl, anchor=west] at (3.4, -2.3) {libre (fragmento desperdiciado)};
\node[paged] at (7.6, -2.3) {};
\node[lbl, anchor=west] at (8.0, -2.3) {bloque paginado reutilizable};
\end{tikzpicture}%
}%
\caption{Gestión del KV cache: esquema clásico (reserva contigua con huecos por fragmentación) frente a PagedAttention (bloques fijos gestionados como memoria virtual). Cada cuadrado representa un fragmento de VRAM del mismo tamaño. Fuente: elaboración propia a partir de~\cite{vllm2023}.}
\label{fig:paged-attention}
\end{figure}

\subsubsection{Batching continuo}

La otra pieza clave del serving moderno es el \textbf{continuous batching} (también llamado \textit{in-flight batching}). Los motores tradicionales procesan peticiones en batches estáticos: toman un grupo de peticiones, les hacen el prefill, generan tokens hasta que todas terminan, y entonces pasan al siguiente batch. Este esquema tiene un problema obvio: una petición corta queda bloqueada esperando a que una petición larga del mismo batch termine, desperdiciando recursos.

En continuous batching, la unidad mínima de trabajo no es ``una petición completa'' sino ``un paso de decode''. En cada iteración, el motor elige qué peticiones vivas tienen un token listo para procesar, las agrupa en un mini-batch y ejecuta un único paso forward para todas. Cuando una petición termina, se elimina sin afectar al resto; cuando llega una nueva, se le hace prefill y se une al pool sin esperar a un batch nuevo. El efecto combinado es que la GPU se mantiene ocupada casi siempre, y las peticiones cortas no pagan la latencia de las largas.

vLLM implementa continuous batching de manera nativa sobre PagedAttention: cada iteración de decode selecciona peticiones vivas, prepara los bloques correspondientes del KV cache y dispatches un kernel CUDA que las procesa en paralelo. Este bucle es lo que verás referido como el \textit{scheduler} en la documentación de vLLM.

\subsubsection{Otras técnicas relacionadas}

Aparte de PagedAttention y continuous batching, hay dos familias de técnicas que conviene mencionar porque son primas cercanas y aparecen frecuentemente en la literatura de serving eficiente.

La primera es \textbf{FlashAttention}~\cite{dao2022flashattention, dao2023flash2}, que optimiza el kernel de atención para reducir los accesos a HBM mediante \textit{tiling} y cómputo en SRAM del chip. FlashAttention es ortogonal a PagedAttention: mientras que PagedAttention ataca la organización del KV cache en memoria, FlashAttention acelera el cálculo del kernel de atención en sí. Las versiones modernas de vLLM lo integran de serie, y ambos beneficios se suman sin interferencias.

La segunda es el \textbf{speculative decoding}~\cite{leviathan2023speculative}, una técnica en la que un modelo ``borrador'' pequeño propone varios tokens por adelantado y el modelo grande los verifica en paralelo. Si los acepta, se gana velocidad porque se generan varios tokens en un solo paso de verificación; si rechaza alguno, se retrocede al punto de divergencia. Es otra forma de reducir la latencia de decode aprovechando hardware infrautilizado durante los pasos de atención del modelo grande. Guarda un parentesco conceptual con la cascada del TFG (dos modelos de distinto tamaño colaboran para servir una petición), pero opera a un nivel distinto, dentro de una misma petición, token a token, y con el modelo grande siempre implicado.

Estas técnicas no son el foco del TFG, pero conviene tenerlas presentes para situar la cascada dentro del panorama más amplio de serving eficiente. En todas ellas, la filosofía es la misma: reducir el coste de servicio sin cambiar significativamente qué se responde.

\subsection{Cascadas, routing y políticas de decisión}

Las partes anteriores de este capítulo han visto cómo entrenar un student que imita al teacher y cómo servir ambos modelos de forma eficiente. Queda la pregunta central del TFG: \textbf{¿cómo decidir, para cada petición, qué modelo debe responder?} Lo que sigue repasa el marco conceptual de cascadas y routing de modelos, y sitúa las tres políticas del TFG (A, B y C) dentro de ese mapa.

\subsubsection{Motivación. No todas las peticiones necesitan el modelo grande}

La observación empírica que da sentido a la cascada es sencilla y muy robusta: en la mayoría de distribuciones de peticiones reales, \textit{la dificultad se concentra en una cola}. Una gran parte de las peticiones son fáciles (preguntas sencillas, aritmética directa, reformulaciones de texto corto) y un modelo pequeño las resuelve correctamente con costes muy inferiores a los del modelo grande. Solo una minoría requiere razonamiento largo, conocimiento de cola o manejo de casos fuera de distribución.

Esta idea no es nueva. En visión por computador o búsqueda de documentos, los \textit{cascade classifiers} llevan décadas aplicándola~\cite{viola2001cascade}. La novedad en LLMs es que el coste por petición es mucho mayor y la diferencia de coste entre modelos de distintos tamaños es enorme: pasar de 14B a 7B puede reducir el coste por token a la mitad o más. Incluso si la política de routing solo acierta en una parte de las peticiones, el ahorro agregado es sustancial.

\subsubsection{Tres familias de políticas}

Aunque en la literatura aparecen muchos esquemas concretos, todos se pueden agrupar en tres familias según \textit{cuándo} deciden qué modelo sirve cada petición.

\textbf{Política A (baseline).} Siempre se usa el teacher. No hay decisión. Sirve como línea base de calidad máxima y coste máximo. Todo el resto se mide contra esta referencia.

\textbf{Política B (cascada forzada).} Se responde primero con el student. Sobre su respuesta se aplica un \textit{quality gate}, es decir, una función barata que decide si la respuesta parece correcta. Si lo es, se acepta; si no, se escala al teacher y se responde con la suya. La ventaja es que no requiere entrenar ningún componente adicional; la desventaja es que las peticiones que acaban escalando pagan tanto el coste del student como el del teacher. Se suele llamar también \textit{inference cascade} o \textit{sequential cascade}~\cite{chen2023frugalgpt}.

\textbf{Política C (cascada por confianza, routing predictivo).} Antes de servir, un componente predictivo mira la petición y estima a qué modelo conviene enviarla. Si predice que el student la resuelve bien, solo se usa el student; si predice que no, se envía directamente al teacher. Se evita así el coste extra del fallback de la Política B, pero a cambio hay que entrenar el predictor y pagar su latencia (aunque sea pequeña). Este es el esquema que usan sistemas como FrugalGPT~\cite{chen2023frugalgpt} y RouteLLM~\cite{dekoninck2024routing}.

\begin{figure}[H]
\centering
\resizebox{\textwidth}{!}{%
\begin{tikzpicture}[
>=Latex,
node distance=6mm and 8mm,
mdl/.style={draw, rounded corners=2pt, align=center, minimum width=22mm, minimum height=9mm, font=\footnotesize, thick},
dec/.style={draw, rounded corners=2pt, align=center, minimum width=18mm, minimum height=9mm, font=\footnotesize, fill=bcrouter!40},
req/.style={draw, rounded corners=2pt, align=center, minimum width=16mm, minimum height=8mm, font=\footnotesize, fill=yellow!30},
outbox/.style={draw, rounded corners=2pt, align=center, minimum width=16mm, minimum height=8mm, font=\footnotesize, fill=bcdata!35}
]
% Políticas en columnas
% --- Col A ---
\node[font=\bfseries\footnotesize] at (0,2.3) {Política A};
\node[req] (ra) at (0,1.2) {petición};
\node[mdl, fill=bcteacher!40] (ta) at (0,0) {Teacher};
\node[outbox] (oa) at (0,-1.2) {respuesta};
\draw[->] (ra) -- (ta);
\draw[->] (ta) -- (oa);

% --- Col B ---
\node[font=\bfseries\footnotesize] at (4.8,2.3) {Política B};
\node[req] (rb) at (4.8,1.2) {petición};
\node[mdl, fill=bcstudent!40] (sb) at (4.8,0) {Student};
\node[dec, minimum width=22mm] (gateb) at (4.8,-1.1) {quality gate};
\node[mdl, fill=bcteacher!40] (tb) at (4.8,-2.4) {Teacher};
\node[outbox] (ob) at (4.8,-3.6) {respuesta};
\draw[->] (rb) -- (sb);
\draw[->] (sb) -- (gateb);
\draw[->] (gateb) -- node[right, font=\scriptsize] {fallback} (tb);
\draw[->] (tb) -- (ob);
\draw[->] (gateb.west) -- ++(-7mm,0) |- node[near start, left, font=\scriptsize] {acepta} (ob.west);

% --- Col C ---
\node[font=\bfseries\footnotesize] at (10,2.3) {Política C};
\node[req] (rc) at (10,1.2) {petición};
\node[dec, minimum width=26mm] (pred) at (10,0) {predictor / router};
\node[mdl, fill=bcstudent!40] (sc) at (8.6,-1.4) {Student};
\node[mdl, fill=bcteacher!40] (tc) at (11.4,-1.4) {Teacher};
\node[outbox] (oc) at (10,-2.8) {respuesta};
\draw[->] (rc) -- (pred);
\draw[->] (pred.west) -- ++(-3mm,0) |- (sc.north) node[midway, left, font=\scriptsize] {fácil};
\draw[->] (pred.east) -- ++(3mm,0) |- (tc.north) node[midway, right, font=\scriptsize] {difícil};
\draw[->] (sc.south) |- (oc.west);
\draw[->] (tc.south) |- (oc.east);
\end{tikzpicture}%
}
\caption{Las tres políticas de cascada comparadas en este TFG. La A siempre usa teacher, la B usa student y escala a teacher solo si el quality gate rechaza, la C decide a priori con un router. Fuente: elaboración propia.}
\label{fig:politicas}
\end{figure}

\subsubsection{Señales usadas por un router}

La Política C requiere un componente que tome la decisión sin ejecutar el modelo grande. Las señales que usa el router pueden venir de tres fuentes distintas. La primera son \textbf{features del prompt}: propiedades que se pueden calcular sin ejecutar ningún modelo, como la longitud, la presencia de números, la densidad de palabras técnicas, el idioma o la similitud con prompts anteriores cacheados. Son baratísimas de calcular pero informan poco sobre la dificultad real de la petición.

La segunda fuente son \textbf{features del student}, que aparecen al hacer inferencia rápida con el modelo pequeño: log-probabilidades de los primeros tokens generados, entropía de la distribución de salida, perplejidad del propio prompt o margen entre el top-1 y el top-2 del último token. Estas señales son mucho más informativas que las anteriores, a costa de tener que pasar la petición por el student para obtenerlas.

La tercera opción son \textbf{features mixtas}, donde un predictor ligero (regresión logística, MLP o \textit{gradient boosting}) combina las señales anteriores para aprender un clasificador de dificultad. Este es el esquema más común cuando se busca un compromiso razonable entre coste del router y capacidad predictiva.

FrugalGPT~\cite{chen2023frugalgpt} propone una estrategia en dos fases: primero una cascada forzada (Política B en nuestros términos), después un router entrenado con las señales acumuladas, y finalmente un \textit{LLM scorer} que decide si aceptar cada respuesta. RouteLLM~\cite{dekoninck2024routing} va más lejos y entrena el router con pares de preferencias (qué modelo era mejor en cada ejemplo) usando aprendizaje preferencial. En ambos casos, la idea común es que el router no necesita ser muy complicado para ser útil: incluso clasificadores lineales sobre señales simples consiguen reducir significativamente el uso del modelo grande.

\subsubsection{Métricas para comparar políticas}

Comparar políticas de cascada no es tan directo como comparar la accuracy de dos modelos. Hay que reportar, como mínimo, tres dimensiones a la vez. La primera es la \textbf{calidad}, medida como accuracy (o \textit{exact match}) en el benchmark de referencia. La segunda es el \textbf{coste}: en el caso de LLMs, la métrica más portable es la fracción de peticiones que acabaron pasando por el teacher, aunque también se suele reportar el coste promedio por petición en términos de tokens generados o GPU-segundo. La tercera es la \textbf{latencia}, especialmente el TTFT y el percentil 95 de la latencia total, porque un router que reduce mucho el uso del teacher pero añade 500~ms de overhead puede no ser útil en una aplicación interactiva.

En el TFG se reportan las tres dimensiones simultáneamente, y se incluye además una comparativa relativa: cuánto porcentaje de calidad se pierde a cambio de cuánto porcentaje de coste se ahorra. Esta es la forma más honesta de caracterizar el compromiso.

\subsubsection{Dónde se sitúa este TFG}

El TFG no pretende inventar un nuevo algoritmo de routing, sino \textbf{caracterizar experimentalmente} las tres políticas bajo el mismo protocolo y en un entorno HPC real. La contribución no es teórica sino metodológica y cuantitativa. Se define un protocolo de evaluación común, reproducible, con warmup, ventanas de medida y repeticiones controladas. Se miden simultáneamente calidad y coste para las tres políticas, y se discute para qué franjas de calidad exigida la cascada es preferible al teacher directo y para cuáles no.

Dentro del mapa del estado del arte, la Política B es un método conocido y bien documentado; la Política C es la variante de routing tipo FrugalGPT/RouteLLM simplificada para un escenario con dos modelos (teacher y student) en lugar de un pool grande. La ambición no es ganar un benchmark sino aportar un estudio honesto y reproducible en MareNostrum~5.

%%%%%%%%%%%%%%%%%%%%%%%%%%%%%%%%%%%%%%%%%%%%%%%%%%%%%%%%%%%%%%%%%
\section{Estado del arte}
\label{sec:estado-del-arte}

Una vez presentados los fundamentos teóricos, conviene situar el trabajo dentro del panorama de investigación reciente. Esta sección revisa las contribuciones más relevantes en las tres líneas que convergen en el TFG (destilación de conocimiento en modelos generativos, serving eficiente de LLMs y routing o cascada entre modelos), con el doble objetivo de acreditar la pertinencia de la propuesta y de delimitar con claridad qué aporta este trabajo respecto a lo que ya existe.

\subsection{Destilación de conocimiento en modelos generativos}

La destilación de conocimiento, formulada originalmente por Hinton et al.~\cite{hinton2015distilling} para modelos de clasificación, ha experimentado una evolución acelerada desde que los LLMs se convirtieron en el paradigma dominante del procesamiento de lenguaje natural. La adaptación al contexto generativo no fue directa: en un clasificador, el teacher produce una distribución sobre un número fijo de clases; en un modelo autoregresivo, la distribución es sobre un vocabulario de decenas de miles de tokens y cambia en cada posición de la secuencia.

Los primeros trabajos en destilar modelos de lenguaje a escala se concentraron en encoders. DistilBERT~\cite{sanh2019distilbert} demostró que un modelo BERT con la mitad de capas, entrenado con una combinación de pérdida de distilación, supervisión clásica y alineamiento de representaciones, conservaba el 97\,\% del rendimiento del teacher en benchmarks de comprensión. Este resultado fue influyente porque estableció un protocolo práctico de destilación que se ha replicado en decenas de variantes posteriores.

La transición a modelos generativos trajo consigo la distinción, ya expuesta en la sección de Knowledge Distillation, entre \emph{word-level KD} y \emph{sequence-level KD}. Kim y Rush~\cite{kim2016sequence} formalizaron esta última para traducción automática y mostraron que entrenar un student con las traducciones del teacher como etiquetas era competitivo con la KD por logits y mucho más sencillo de implementar. Este enfoque fue el precursor directo de la oleada de modelos destilados que apareció en 2023--2024: Alpaca~\cite{taori2023alpaca} (student de LLaMA entrenado sobre salidas de GPT-3.5), Vicuna (sobre conversaciones de ShareGPT), Orca (que además incorpora cadenas de razonamiento del teacher), y la familia Phi de Microsoft (que combina destilación con curación cuidadosa de datos sintéticos).

En paralelo, Agarwal et al.~\cite{agarwal2024gkd} propusieron \textit{Generalized Knowledge Distillation} (GKD), un esquema on-policy donde el student genera y el teacher evalúa. GKD aborda el problema de \textit{train--test mismatch} propio de la destilación off-policy y reporta mejoras consistentes en tareas de generación abierta, aunque a costa de un presupuesto de cómputo significativamente mayor. La tensión entre off-policy (barato, reproducible, reutilizable) y on-policy (más fiel, más caro) es un eje central en la literatura actual y ha condicionado directamente la elección de este TFG, que opta por sequence-level off-policy por razones de presupuesto.

Un resultado reciente especialmente relevante es el de la familia Qwen 2.5~\cite{qwen2}, que incluye modelos desde 0,5B hasta 72B parámetros dentro de la misma arquitectura. Los modelos pequeños de esta familia se benefician de destilación durante su pipeline de entrenamiento, lo que los convierte en puntos de partida ya parcialmente destilados. Este hecho es importante para el TFG: el student de 1,5B que se usa como base no parte de cero en términos de capacidad, sino de un modelo que ya ha sido expuesto a señales de modelos mayores durante su preentrenamiento.

\subsection{Serving eficiente de modelos de lenguaje}

El serving de LLMs ha sido un área de innovación rápida impulsada por la necesidad industrial de reducir el coste por token generado. PagedAttention~\cite{vllm2023}, incorporado en vLLM, fue el primer salto cualitativo al aplicar la idea de paginación de memoria virtual al KV cache, multiplicando el throughput sin alterar la calidad. Desde entonces, el ecosistema se ha diversificado con motores como TensorRT-LLM (NVIDIA), SGLang, y TGI (HuggingFace), cada uno con optimizaciones especializadas.

FlashAttention~\cite{dao2022flashattention, dao2023flash2} atacó el cuello de botella de la atención por otro flanco: reorganizando el cómputo para minimizar los accesos a HBM mediante \textit{tiling} y fusión de operaciones. La combinación de FlashAttention (kernel rápido) con PagedAttention (gestión eficiente de memoria) es hoy el estándar de facto en serving de LLMs y es exactamente la configuración que usa el TFG a través de vLLM.

El \textit{speculative decoding}~\cite{leviathan2023speculative} introdujo un paradigma distinto, en el que un modelo borrador propone tokens que el modelo principal verifica en paralelo. Aunque guarda un parentesco conceptual con la cascada del TFG (dos modelos de distinto tamaño colaboran), opera a una escala diferente (dentro de una misma petición, token a token) y con el modelo grande siempre involucrado. No obstante, es técnicamente compatible con el enfoque de cascada, de modo que un student destilado podría servir simultáneamente como modelo borrador para speculative decoding cuando la petición se enruta al teacher.

Más recientemente, la investigación se ha orientado hacia la desagregación de prefill y decode en GPUs separadas, la cuantización dinámica del KV cache, y el uso de caches jerárquicas que persisten entre peticiones. Estas líneas son complementarias a la cascada y podrían integrarse en futuras iteraciones del sistema.

\subsection{Routing y cascadas de modelos}

La idea de asignar dinámicamente peticiones a modelos de distinto coste es anterior a los LLMs. Los \textit{cascade classifiers} de Viola y Jones~\cite{viola2001cascade} ya aplicaban este principio en detección de caras, con un clasificador rápido filtrando la mayoría de ventanas negativas antes de invocar al costoso. En el contexto de LLMs, FrugalGPT~\cite{chen2023frugalgpt} fue uno de los primeros trabajos en proponer explícitamente una cascada de modelos de lenguaje, combinando un modelo barato con un \textit{quality scorer} y un modelo caro, y demostrando ahorros de coste de hasta un 98\,\% con pérdidas de calidad inferiores al 1\,\% en ciertos benchmarks.

RouteLLM~\cite{dekoninck2024routing} formalizó la idea con un router entrenado mediante aprendizaje preferencial sobre pares de respuestas, y demostró que incluso un clasificador simple puede desviar una fracción significativa del tráfico al modelo pequeño sin pérdida apreciable de calidad percibida. El trabajo de Ong et al.~\cite{ong2025routellm} propuso un marco unificado que engloba tanto el routing (decisión ex-ante sobre a qué modelo enviar) como el cascading (escalado secuencial después de una primera respuesta), y definió una función de utilidad
\[
u(x) \;=\; q(x) \;-\; \lambda\, c(x),
\]
donde $q(x)$ es la calidad esperada, $c(x)$ el coste, y $\lambda$ un parámetro que controla el compromiso. Esta formulación es la base conceptual de la Política C del TFG.

La propuesta de este TFG se diferencia de los trabajos anteriores en un aspecto específico, la redefinición del coste $c(x)$ como \textbf{coste de servicio}, incorporando el estado dinámico del sistema (ocupación de GPUs, tamaño de batch, longitud de cola, KV cache residual) como variable del problema. Mientras que FrugalGPT y RouteLLM definen $c(x)$ como un coste estático asociado al modelo (precio por API call, o latencia media), la extensión propuesta redefine la utilidad como $u(x,z) = q(x) - \lambda\,c(x,z)$, donde $z$ captura el estado del sistema de serving en el instante de la decisión. Esta formulación busca demostrar que un router consciente del estado del sistema puede mantener la misma calidad y latencia con mejor throughput útil y menor consumo energético.

\subsection{Posicionamiento del TFG}

Con este panorama, el TFG se sitúa en la intersección de las tres líneas: utiliza la destilación de conocimiento para producir un student eficiente, despliega teacher y student sobre un motor de serving moderno en un entorno HPC real, y evalúa políticas de cascada que combinan ambos modelos. La contribución no pretende ser teórica en ninguna de las tres áreas, sino metodológica e integradora: construir un pipeline completo, reproducible y medible que conecte las tres piezas en MareNostrum~5, y caracterizar experimentalmente el compromiso calidad-coste-servicio bajo condiciones realistas. La extensión del coste a coste de servicio, sensible al estado del sistema, constituye la novedad conceptual más relevante del proyecto.

\subsection{Entorno HPC en MareNostrum~5 y SLURM}

La última pieza del marco teórico es el entorno donde todo esto se ejecuta. Este apartado describe, con el nivel de detalle necesario para entender el diseño del pipeline, cómo es MareNostrum~5 y qué implicaciones tiene sobre el proyecto.

\subsubsection{MareNostrum~5 a vista de pájaro}

MareNostrum~5 es el supercomputador del Barcelona Supercomputing Center (BSC), uno de los centros de referencia de computación de altas prestaciones en Europa~\cite{mn5}. Incluye particiones generales de CPU y particiones de aceleradores GPU. La relevante para este TFG es la partición de GPUs NVIDIA Hopper (H100): cada nodo combina CPUs de altas prestaciones con múltiples GPUs H100 de 64~GB de VRAM conectadas por NVLink, y a nivel de clúster los nodos se comunican por una red de alta velocidad tipo InfiniBand.

Desde el punto de vista del usuario, el clúster es un sistema Linux al que se entra por SSH y donde no se ejecutan comandos directamente en las GPUs: se envían \textit{jobs} al gestor de colas. Cuando el gestor encuentra recursos libres, lanza el job, lo ejecuta y devuelve los resultados. Esto introduce una latencia entre ``quiero ejecutar algo'' y ``se está ejecutando'' que puede ir desde segundos hasta horas, según la ocupación del clúster.

\subsubsection{SLURM como gestor de colas}

El gestor de colas es \textbf{SLURM}~\cite{slurm}. SLURM es el estándar de facto en la mayoría de supercomputadores académicos y ofrece todo lo esencial: solicitar recursos (nodos, CPUs, GPUs, memoria, tiempo), controlar la cola, monitorizar jobs, limitar el uso por usuario o grupo y registrar consumo de recursos (incluido \textit{energía} vía \texttt{sacct}). Sus comandos más habituales son \texttt{sbatch} para enviar un job, \texttt{squeue} para ver el estado de la cola, \texttt{scancel} para cancelar, \texttt{sacct} para ver el histórico y \texttt{salloc}/\texttt{srun} para sesiones interactivas.

Cada job se describe con un script de bash anotado con directivas \texttt{\#SBATCH}. Para este TFG, los scripts de experimentación solicitan una GPU H100 completa, una cantidad fija de CPUs y un tiempo límite ajustado al tipo de experimento (desde minutos para un benchmark hasta horas para una destilación). La política de la cola se rige por la cuota asignada al grupo de investigación Data Centric Computing.

\subsubsection{Implicaciones para el proyecto}

El entorno HPC no es solo ``ordenadores más grandes''; cambia cómo hay que diseñar los experimentos. Para este TFG, las implicaciones más importantes son las siguientes.

La primera es que \textbf{todo experimento es un job}. No se puede mantener una sesión interactiva abierta con GPU durante 19 semanas. Cada experimento se describe como un script SLURM reproducible, y cada ejecución queda identificada por su JobID. Esto obliga a tener automatización desde el primer día, pero a cambio proporciona trazabilidad completa de cada ejecución.

La segunda es que \textbf{la ocupación de la cola introduce varianza temporal}. El mismo job puede arrancar a los 2 minutos o a las 6 horas, dependiendo de la carga del sistema. Esto afecta tanto a la planificación como a las comparaciones. Si se quieren comparar dos configuraciones, es necesario asegurarse de que se ejecutan en condiciones equivalentes (mismo tipo de nodo, misma configuración de red, idealmente dentro de una ventana temporal cercana).

La tercera es que \textbf{la infraestructura de red del clúster impone sus propias restricciones}. Servir un LLM con una API HTTP dentro de un nodo del clúster es posible, pero el acceso desde fuera del nodo o desde fuera del clúster requiere túneles SSH o configuraciones específicas. Por ello, el pipeline de evaluación se diseña para que el cliente se ejecute dentro del mismo nodo que el servidor, evitando atravesar cualquier frontera de red.

La cuarta es que \textbf{la energía se puede medir}. SLURM registra \texttt{ConsumedEnergy} por job, lo que permite comparar el coste energético real de las distintas políticas de cascada. En la sección de sostenibilidad ya se avanzó que este valor se registra y se publica junto con los resultados experimentales.

Con esto se cierra el marco teórico. Quedan establecidas todas las piezas necesarias, entre ellas los fundamentos de redes neuronales y modelos de lenguaje, la arquitectura transformer y el coste de su inferencia, el \textit{fine-tuning} y el PEFT, la destilación de conocimiento, el serving moderno con vLLM, las cascadas y el routing entre modelos, y el entorno HPC donde se ejecuta todo. En la Parte~\ref{part:diseno-sistema} se utiliza esta base para describir el diseño del sistema.

%%%%%%%%%%%%%%%%%%%%%%%%%%%%%%%%%%%%%%%%%%%%%%%%%%%%%%%%%%%%%%%%%
